\chapter{I/O、异步与背压}

计算系统不只做 CPU 运算。磁盘、网络、管道、日志和用户请求都会引入 I/O。I/O 的延迟和吞吐特征与 CPU 完全不同，因此运行时必须区分计算等待和 I/O 等待。

\section{从阻塞 I/O 到异步流水线}

同步 I/O 调用会让当前线程等待结果。少量线程处理少量 I/O 时简单可靠；大量并发连接或慢设备场景下，同步等待会浪费线程。异步 I/O 把提交和完成分离，让线程可以处理其他工作。

事件循环、回调、future、协程和 \code{io_uring} 都是组织异步 I/O 的方式。选择哪种方式，要看系统需要的并发数、延迟要求、代码复杂度和平台支持。

异步 I/O 的核心不是“换一种语法”，而是把等待从工作线程上移走。同步调用让当前线程停在系统调用或库调用里；异步模型把请求提交出去，完成时再处理结果。这样可以用较少线程管理大量等待。但异步也会引入状态机、取消、超时、错误传播和资源生命周期问题。若并发度很低，同步 I/O 可能更简单可靠。

\section{背压、队列和批处理}

背压是系统告诉上游“我处理不过来”的机制。没有背压，队列会无限增长，内存耗尽，延迟失控。背压可以通过 bounded queue、令牌桶、限流、拒绝请求、降级和批处理实现。

并发系统里 bounded queue 的意义不只是节省内存，更是把过载变成可观察、可控制的状态。阻塞、返回错误或丢弃，都比默默无限排队更可管理。

背压必须一路传到真正的生产者。只在中间队列限长，但上游仍然无限接收请求，过载会换一个地方堆积。一个服务端如果 worker 队列满了，可以选择拒绝新请求、降低读取速度、返回重试、丢弃低优先级任务或触发降级。选择哪种方式取决于业务语义，但不能没有选择。

\topic{批处理}

批处理用一次调度或一次 I/O 处理多个元素，摊薄固定成本。网络发送、日志写入、数据库提交和向量计算都常用批处理。批太小浪费固定成本，批太大增加延迟和缓存压力。

批处理要控制两个阈值：数量阈值和时间阈值。只按数量凑批，在低流量时可能让请求等太久；只按时间刷批，在高流量时可能批太小。很多系统使用“达到 N 条或等待 T 时间就发送”的策略。批处理还要考虑失败语义：一批里部分成功时如何重试，是否会重复执行，是否需要幂等 ID。

\topic{CPU 与 I/O 的边界}

I/O 密集任务不应占满计算线程池。常见设计是分离 I/O 线程、计算线程和后台维护线程，或者使用异步运行时统一调度但明确资源限制。否则慢 I/O 会阻塞计算任务，计算任务也会延迟 I/O 完成处理。

\topic{队列是系统的压力表}

队列长度、入队速率、出队速率和等待时间，是理解系统压力的基本指标。队列持续增长说明下游处理能力不足或上游没有背压；队列经常为空说明下游可能在等上游；队列长度锯齿状变化可能是批处理或周期性调度造成的。只看 CPU 使用率可能误判，因为系统可能 CPU 很空但卡在 I/O，也可能 CPU 很满但队列仍然增长。

有界队列实验要记录三种结果：生产者是否阻塞，消费者是否空闲，丢弃或拒绝是否发生。这样才能判断背压策略是否真的控制住过载。

\topic{I/O 和 CPU 的成本模型不同}

CPU 计算的基本问题是怎样让核心持续退休有用指令，I/O 的基本问题是怎样管理等待、排队和设备吞吐。一次磁盘读取、网络请求或管道写入可能比一条 CPU 指令慢很多数量级。若用 CPU-bound 的思维处理 I/O，很容易开太多线程、阻塞 worker、隐藏队列积压，最后系统看起来 CPU 不满但请求延迟很高。

I/O 还有两个重要特征。第一，吞吐和延迟经常需要批处理折中。单条写入延迟低但吞吐差，大批写入吞吐高但单条等待更久。第二，错误和取消是常态。磁盘可能返回短读，网络可能超时，远端可能重置连接，写入可能部分成功。高质量运行时不能只写成功路径。

Compute Systems Engine 后续的分布式 shuffle 会大量使用 I/O 思维。map 端产生中间数据，缓冲、批量写入、按分区发送、等待 reduce 端接收、处理重试和背压。如果单机阶段没有建立 I/O 模型，分布式阶段就会写成“把数据发过去”这种空话。

\topic{同步阻塞的适用边界}

同步阻塞 I/O 并不是低级。对于简单工具、并发度低的批处理程序、启动脚本或一次性文件读取，同步 I/O 清晰可靠。问题出在大量连接或大量慢设备等待时，每个等待都占用一个线程。线程有栈、调度状态、上下文切换成本和内存占用；阻塞线程太多，系统会被调度和内存压力拖垮。

因此选择同步还是异步，不是看语法潮流，而是看并发等待数量和生命周期复杂度。若程序同时等待几十个文件读取，同步线程池可能足够；若服务同时维持几十万连接，事件循环或异步 I/O 更合适。若任务既有 CPU 密集计算又有慢 I/O，至少要把计算线程池和 I/O 等待分开，避免慢 I/O 占满计算 worker。

\topic{事件循环的核心}

事件循环把“等待多个事件”集中到少量线程中。它通常维护一组已注册的 fd、定时器和任务队列，等待内核通知哪些事件 ready，然后执行对应回调。Linux 上常见机制包括 \code{epoll} 和 \code{io_uring}。事件循环不等于更快的 CPU，它只是避免每个等待都占用一个线程。

事件循环的难点是回调不能长时间占用 loop 线程。若在事件循环线程里做 CPU 密集解析、压缩或排序，其他 I/O 完成事件会被延迟处理。常见架构是事件循环负责接收和发送，CPU 工作提交到计算线程池，完成后再把结果交回事件循环。这样系统同时需要 I/O 队列、计算队列和背压策略。

\topic{\texorpdfstring{\code{io_uring}}{io-uring} 的直觉}

\code{io_uring} 的核心直觉是提交队列和完成队列。用户态把 I/O 请求放入提交队列，内核处理后把结果放入完成队列。它减少某些系统调用和数据结构转换成本，也支持更丰富的异步操作。但它不是魔法：设备仍有真实延迟和带宽，队列仍可能满，请求仍可能失败，完成处理仍需要 CPU。

学习 \code{io_uring} 时，不要一开始追求所有高级特性。先理解三件事：提交多少 in-flight 请求，完成事件如何被消费，失败和取消如何传播。过多 in-flight 请求会增加队列延迟和内存占用；完成处理太慢会让完成队列积压；取消语义不清会导致请求完成后访问已释放对象。

\section{超时、取消和错误传播}

异步系统必须处理超时和取消。一个请求提交出去后，上层可能已经超时返回，用户对象可能准备释放，但底层 I/O 仍可能稍后完成。若完成回调访问已销毁对象，就会出现 use-after-free。正确设计通常需要请求上下文对象、引用计数或所有权转移，确保完成路径和取消路径共享同一生命周期协议。

超时也不是简单地“过了时间就删掉”。如果底层操作不能立即取消，超时只表示上层不再等待，底层完成事件仍要被回收和丢弃。若请求已经发送到远端，远端可能仍执行，重试可能造成重复提交。因此超时、取消和幂等会在分布式章节重新出现。单机异步 I/O 是理解分布式失败语义的前置训练。

\topic{背压的几种策略}

背压不是只有阻塞。常见策略包括阻塞等待、立即失败、带超时等待、丢弃低优先级任务、降级处理、批量合并、采样、令牌桶限流和动态降低读取速度。每种策略都有业务语义。日志采样可以丢一部分低优先级日志；支付请求不能随便丢；监控指标可以合并；用户请求可能返回 429 或重试建议。

有界队列是背压的最小机制。队列满时，生产者必须面对现实：等、失败、丢弃或降级。无限队列只是把这个选择推迟到 OOM 或尾延迟爆炸。第二册里的 bounded queue 不只是并发练习，它是流控思想的基础。

\begin{lstlisting}[language=C++,caption={带超时的提交语义草图}]
enum class SubmitResult : std::int32_t {
    Accepted,
    TimedOut,
    Closed,
};

struct SubmitPolicy {
    std::int32_t timeout_milliseconds = 0;
    bool allow_drop = false;
};
\end{lstlisting}

这个片段没有实现队列，而是表达接口语义。调用者需要知道任务是否被接受、是否超时、队列是否关闭。若接口只返回 void，调用者无法区分提交成功和任务丢失。系统工程里，返回值设计就是背压协议的一部分。

\topic{批处理的延迟预算}

批处理必须有延迟预算。假设日志写入器每满 4096 条刷一次磁盘，如果流量很低，最后几条日志可能等很久。因此通常设置两个阈值：数量达到 N 或时间达到 T 就刷。N 控制吞吐，T 控制尾延迟。选择 N 和 T 要看业务：监控指标可以延迟几秒，交互请求响应不能等太久。

批处理还要考虑内存和缓存。批太大，可能超过缓存，处理时局部性下降；批太小，系统调用和队列成本占比高。网络批处理还要考虑 MTU、拥塞控制和远端处理能力。数据库提交还要考虑事务语义和 fsync 成本。没有一个“最佳批大小”，只有工作负载下的实验矩阵。

\topic{异步流水线}

一个典型异步流水线可以分成读取、解析、计算、写出四段。读取是 I/O 密集，解析可能 CPU 密集，计算可能需要线程池，写出又是 I/O 密集。段与段之间用有界队列连接。每个队列的长度、入队速率和出队速率告诉我们瓶颈在哪。

\begin{lstlisting}[numbers=none]
read stage -> parse queue -> compute workers -> write queue -> output stage
    I/O          bounded          CPU              bounded       I/O
\end{lstlisting}

若 parse queue 持续增长，说明解析或计算跟不上读取；若 write queue 持续增长，说明输出慢；若所有队列为空而 CPU 很空，说明上游读取慢；若 CPU 满且队列也增长，说明计算是瓶颈。队列指标让系统不再是黑箱。

\topic{优先级和公平性}

真实系统通常有不同优先级任务。高优先级请求、后台压缩、指标上报、日志写入不应完全同等对待。单个 FIFO 队列可能让大量低优先级任务阻塞高优先级请求。可以使用多队列、优先级队列、配额、公平调度或令牌桶隔离。

优先级也会带来饥饿风险。若高优先级持续不断，低优先级任务可能永远不执行。公平性策略要说明最低保障、老化机制或后台窗口。系统设计经常是在延迟、吞吐、公平和实现复杂度之间取舍。

\topic{I/O 错误和重试}

I/O 错误必须分类。临时错误可以重试，例如网络超时、EAGAIN、远端暂时不可用；永久错误不应盲目重试，例如权限错误、格式错误、目标不存在。重试要有预算、退避和幂等保障。无限立即重试会放大故障，形成重试风暴。

写入操作尤其要小心部分成功。一次批量写可能只写了一部分；一次远程请求可能已经被对方执行但响应丢失；一次文件写入可能进入 page cache 但还没持久化。系统必须定义提交点：什么时候认为任务完成，什么时候可以重试，重复执行是否安全。分布式章节会把这套语义扩展到幂等 ID 和检查点。

\section{Linux 观察、可靠写入和项目落地}

Linux I/O 观察可以从几个层次开始。系统调用层可用 \cmd{strace} 观察 \code{read}/\code{write}/\code{epoll}/\code{io_uring} 相关调用，但它会带来明显开销。性能层可以用 \cmd{perf} 看 CPU 热点和调度，使用 \cmd{iostat}、\cmd{pidstat} 或 \cmd{ss} 观察设备和网络状态。内核源码阅读可以从 VFS、page cache、block layer、epoll 或 \code{io_uring} 子系统选择一个窄问题进入。

不要写“Linux I/O 很复杂”这种泛泛结论。应围绕现象阅读：为什么 mmap 首次访问触发 page fault，为什么 epoll 能等待多个 fd，为什么写文件返回不等于持久化，为什么 \code{io_uring} 有提交队列和完成队列。读源码要和用户态实验绑定。


\topic{从一次具体事故开始}

reduce 生成输出比磁盘写得快，CPU 看起来很空，内存却持续上涨，最终 checkpoint 还留下半写文件。这不是抽象练习，而是后续所有机制的入口。计算线程、I/O 等待、完成事件和持久化提交被混在一起，系统没有把下游容量反馈给上游。如果这一层只写成术语列表，读者会知道很多名词，却不知道面对这类现象时先固定什么、再观察什么、最后改什么。

本章把运行问题固定为读取输入文件、写 shuffle block、提交 checkpoint。这件事的价值是让所有推导都落在同一条数据流上：输入如何进入系统，状态在哪里改变，输出何时对外可见，失败发生后哪些东西还能相信。主失败可以压缩成一句话：执行和完成分离后，生命周期、背压、超时和持久化边界变得比函数调用复杂。后面的每个机制都要回到这个失败，而不是各讲各的。

\begin{lstlisting}[numbers=none]
compute block -> bounded write queue
              -> submit io -> completion event
              -> fsync temp -> rename -> manifest
\end{lstlisting}

阅读本章时要先画这条链路，而不是先背概念。链路中每个箭头都表示一次边界跨越：从文件到内存，从线程到队列，从局部状态到共享状态，从临时结果到提交结果。边界越多，隐含假设越多；隐含假设越多，优化时越容易把正确性、性能和恢复搅在一起。

本章的目标不是给出一个万能框架，而是训练一种判断顺序：先写 reference，再暴露朴素版本，再沿着失败路径引入机制，最后用实验和审查表确定结论边界。如果一个优化不能说明它改变了哪条边界，就先不要把它合进贯穿项目。

\topic{朴素方案为什么会被写出来}

第一版计算线程直接 write；慢了以后改成后台线程无界缓存；checkpoint 直接覆盖旧文件。这个写法并不是一开始就荒谬。它符合小程序直觉：对象少、状态近、调用栈清楚、失败罕见，所有问题都能在一个调试会话里看见。经典教材写到这里不能直接嘲笑朴素方案，因为读者需要先理解它在哪些条件下成立。

朴素方案的真正问题，是它把多个边界藏了起来。控制状态和业务数据混在一起，临时结果和提交结果混在一起，测量边界和语义边界混在一起，线程资源和任务数量混在一起。代码短只是表面现象；代价是没有地方记录不变量，也没有地方插入观测点。

把这个方案放大到读取输入文件、写 shuffle block、提交 checkpoint，第一个变化是输入规模让偶发路径变成常态，第二个变化是并发让顺序假设失效，第三个变化是系统资源让等待和数据移动变成主成本。所以改进不是在原函数里继续加分支，而是把边界显式化。

教材里的朴素版本必须保留。它一方面作为 correctness oracle 的对照，另一方面作为解释机制的反面样本。没有朴素版本，最终设计就像凭空出现；没有反例，最终设计又会被误解成永远正确。

\topic{失败路径一：结果错误不是性能问题}

无界写队列让内存承担所有下游慢速，超时后底层请求仍可能完成，buffer 生命周期不清。这类问题通常不会稳定复现。小输入没有触发，单线程没有触发，甚至同一批输入换一个调度顺序才触发。因此第一反应不能是调线程数，也不能是继续优化热循环，而是先把语义锚点拿出来。

排查时按四步走。第一步固定输入和随机种子；第二步运行串行 reference，保存可比较摘要；第三步让优化版本输出同样的摘要；第四步定位第一个不同的边界。摘要不能只写总数，还要包含错误分类、分片身份、attempt 身份、输出 checksum 和阶段状态。

\begin{lstlisting}[numbers=none]
t0 run reference and save checksum
t1 run naive implementation on the same input
t2 compare record count, error count, key count, output checksum
t3 stop at the first boundary that diverges
\end{lstlisting}

这里的关键是“第一个不同的边界”。如果 reference 与优化版本在解析阶段已经不同，就不要继续讨论 cache、SIMD、线程池或分布式提交。系统分析要避免把下游症状当成根因。

\topic{失败路径二：吞吐提升可能掩盖资源失控}

write 返回不等于可恢复提交，进程崩溃可能留下 page cache 中未持久化的数据或半截 checkpoint。这类失败常常伴随一个迷惑性信号：平均吞吐看起来变好，尾延迟、内存峰值、恢复时间或输出可信度却变坏。高质量教材必须把这类反例写进正文，否则读者会误以为所有优化只要更快就是更好。

资源失控要按队列、内存、文件、线程和网络五类检查。队列看水位和等待，内存看峰值和分配次数，文件看临时产物和持久化点，线程看 runnable 与 blocked，网络或消息看 in-flight 与重试。单个总耗时无法解释这些状态。

\begin{lstlisting}[numbers=none]
observe:
  queue_depth, producer_wait_ns, consumer_wait_ns
  resident_bytes, allocations, temporary_files
  runnable_threads, blocked_threads, retry_count
  committed_outputs, stale_outputs, checksum_errors
\end{lstlisting}

如果某个指标没有采集，就在报告里明确写成未验证风险。不要把没有观测到的路径写成不存在，也不要用当前机器的一次稳定运行去替代边界证明。

\topic{把机制落成状态和接口}

本章把 I/O 视为状态机：提交、排队、完成、取消、错误分类、fsync、rename 和 manifest 都要明确。状态和接口是机制进入工程的地方。概念如果不能落成字段、状态、函数边界、错误码或实验指标，就很难在代码审查里被检查。

\begin{lstlisting}[language=C++,caption={IoBackpressureState 的最小教学结构}]
struct WriteRequest {
    std::uint64_t block_id = 0;
    std::uint64_t bytes = 0;
    std::uint64_t checksum = 0;
    std::uint64_t deadline_ns = 0;
};

struct WriteCompletion {
    std::uint64_t block_id = 0;
    std::int32_t error_code = 0;
    std::uint64_t completed_bytes = 0;
};
\end{lstlisting}

这段结构不是要替代真实项目代码，而是把本章最容易被忽略的边界写出来。字段名应回答一个具体问题：身份是什么，状态属于谁，哪次 attempt 有效，哪些字节已经被处理，哪个 checksum 能证明输出没有被误读。

接口设计还有一个原则：控制面字段不要被热路径随手修改，数据面批量结构不要夹带恢复协议。前者让状态机可审查，后者让硬件看到规则的数据流。若两者混在一起，代码会同时难以优化和难以恢复。

\topic{机制推导：阻塞 I/O}

先看一个具体问题：read 阻塞时 worker 到底在等什么。在读取输入文件、写 shuffle block、提交 checkpoint中，这个问题会沿着输入、执行、同步、I/O 或提交边界继续传播，不会停留在一个局部函数里。

朴素写法通常是：在计算线程中直接读写文件。它吸引人的地方是短、直观、容易在小样本上通过测试；危险也在这里，因为小样本会替它隐藏资源、调度、数据分布和失败路径。

放大以后会出现的信号是：线程被 I/O 等待占住，计算资源和等待资源混在一起。这个信号未必是崩溃，也可能是吞吐平台期、p99 抖动、结果偶尔变化、队列持续变长、重启后状态不一致，或者某个分片长期拖尾。

此时引入阻塞 I/O，不是为了增加术语，而是为了建立一个可检查模型：阻塞 I/O 简单但会把等待带入执行线程。模型必须同时解释正确性和成本。只解释为什么快，不解释为什么可信，不能进入贯穿项目；只解释为什么正确，却完全无视缓存、调度、I/O 或网络，也不符合第二册目标。

把模型写成不变量时，要问四个问题：哪些状态只能由一个拥有者修改，哪些结果允许重试，哪些数据可以批量搬运，哪些错误必须外显。围绕阻塞 I/O的代码审查，就从这四个问题开始，而不是从实现看起来是否高级开始。

观察方式应围绕假设设计：记录 blocked time、CPU 利用率和队列长度。实验要固定输入和环境，保留 reference，一次只改变一个变量。如果某项工具在当前手机或 Linux 环境没有权限，就写出预期命令、缺失事件和结论限制。

最后要写边界：对小工具阻塞 I/O 可能足够，不必过度异步。边界不是附录里的保守声明，而是正文的一部分。读者应该知道什么时候使用它，什么时候退回简单方案，什么时候需要换一个尺度重新测。

本节验收可以用一句话判断：能否从朴素版本推导到改进版本，并能说清新增字段、新增状态或新增实验到底保护了什么。如果只能给最终代码，却解释不了这条推导链，说明机制还没有真正学会。

\topic{机制推导：异步完成}

先看一个具体问题：提交请求后，结果何时回来，谁拥有 buffer。在读取输入文件、写 shuffle block、提交 checkpoint中，这个问题会沿着输入、执行、同步、I/O 或提交边界继续传播，不会停留在一个局部函数里。

朴素写法通常是：把异步调用当成立刻完成。它吸引人的地方是短、直观、容易在小样本上通过测试；危险也在这里，因为小样本会替它隐藏资源、调度、数据分布和失败路径。

放大以后会出现的信号是：buffer 生命周期、错误返回和取消都不清楚。这个信号未必是崩溃，也可能是吞吐平台期、p99 抖动、结果偶尔变化、队列持续变长、重启后状态不一致，或者某个分片长期拖尾。

此时引入异步完成，不是为了增加术语，而是为了建立一个可检查模型：异步 I/O 把提交和完成拆开，需要 completion 事件和所有权规则。模型必须同时解释正确性和成本。只解释为什么快，不解释为什么可信，不能进入贯穿项目；只解释为什么正确，却完全无视缓存、调度、I/O 或网络，也不符合第二册目标。

把模型写成不变量时，要问四个问题：哪些状态只能由一个拥有者修改，哪些结果允许重试，哪些数据可以批量搬运，哪些错误必须外显。围绕异步完成的代码审查，就从这四个问题开始，而不是从实现看起来是否高级开始。

观察方式应围绕假设设计：测试完成前释放 buffer 和取消场景。实验要固定输入和环境，保留 reference，一次只改变一个变量。如果某项工具在当前手机或 Linux 环境没有权限，就写出预期命令、缺失事件和结论限制。

最后要写边界：异步增加状态机复杂度，必须有清晰封装。边界不是附录里的保守声明，而是正文的一部分。读者应该知道什么时候使用它，什么时候退回简单方案，什么时候需要换一个尺度重新测。

本节验收可以用一句话判断：能否从朴素版本推导到改进版本，并能说清新增字段、新增状态或新增实验到底保护了什么。如果只能给最终代码，却解释不了这条推导链，说明机制还没有真正学会。

\topic{机制推导：背压}

先看一个具体问题：写出速度慢于计算速度时，上游应该怎么办。在读取输入文件、写 shuffle block、提交 checkpoint中，这个问题会沿着输入、执行、同步、I/O 或提交边界继续传播，不会停留在一个局部函数里。

朴素写法通常是：继续产生输出并排队。它吸引人的地方是短、直观、容易在小样本上通过测试；危险也在这里，因为小样本会替它隐藏资源、调度、数据分布和失败路径。

放大以后会出现的信号是：内存增长、延迟增长，最终 OOM 或超时。这个信号未必是崩溃，也可能是吞吐平台期、p99 抖动、结果偶尔变化、队列持续变长、重启后状态不一致，或者某个分片长期拖尾。

此时引入背压，不是为了增加术语，而是为了建立一个可检查模型：背压把下游容量反馈到上游，让生产速度受限制。模型必须同时解释正确性和成本。只解释为什么快，不解释为什么可信，不能进入贯穿项目；只解释为什么正确，却完全无视缓存、调度、I/O 或网络，也不符合第二册目标。

把模型写成不变量时，要问四个问题：哪些状态只能由一个拥有者修改，哪些结果允许重试，哪些数据可以批量搬运，哪些错误必须外显。围绕背压的代码审查，就从这四个问题开始，而不是从实现看起来是否高级开始。

观察方式应围绕假设设计：记录队列水位、生产者等待和丢弃策略。实验要固定输入和环境，保留 reference，一次只改变一个变量。如果某项工具在当前手机或 Linux 环境没有权限，就写出预期命令、缺失事件和结论限制。

最后要写边界：背压策略是业务语义，不能只有技术默认。边界不是附录里的保守声明，而是正文的一部分。读者应该知道什么时候使用它，什么时候退回简单方案，什么时候需要换一个尺度重新测。

本节验收可以用一句话判断：能否从朴素版本推导到改进版本，并能说清新增字段、新增状态或新增实验到底保护了什么。如果只能给最终代码，却解释不了这条推导链，说明机制还没有真正学会。

\topic{机制推导：批处理}

先看一个具体问题：为什么每条记录一次 write 会慢。在读取输入文件、写 shuffle block、提交 checkpoint中，这个问题会沿着输入、执行、同步、I/O 或提交边界继续传播，不会停留在一个局部函数里。

朴素写法通常是：记录一来就写出。它吸引人的地方是短、直观、容易在小样本上通过测试；危险也在这里，因为小样本会替它隐藏资源、调度、数据分布和失败路径。

放大以后会出现的信号是：系统调用、锁和设备提交成本被放大。这个信号未必是崩溃，也可能是吞吐平台期、p99 抖动、结果偶尔变化、队列持续变长、重启后状态不一致，或者某个分片长期拖尾。

此时引入批处理，不是为了增加术语，而是为了建立一个可检查模型：批处理让一次昂贵操作服务多条记录。模型必须同时解释正确性和成本。只解释为什么快，不解释为什么可信，不能进入贯穿项目；只解释为什么正确，却完全无视缓存、调度、I/O 或网络，也不符合第二册目标。

把模型写成不变量时，要问四个问题：哪些状态只能由一个拥有者修改，哪些结果允许重试，哪些数据可以批量搬运，哪些错误必须外显。围绕批处理的代码审查，就从这四个问题开始，而不是从实现看起来是否高级开始。

观察方式应围绕假设设计：比较批大小、延迟分位数和吞吐。实验要固定输入和环境，保留 reference，一次只改变一个变量。如果某项工具在当前手机或 Linux 环境没有权限，就写出预期命令、缺失事件和结论限制。

最后要写边界：批太大增加尾延迟和崩溃后重做成本。边界不是附录里的保守声明，而是正文的一部分。读者应该知道什么时候使用它，什么时候退回简单方案，什么时候需要换一个尺度重新测。

本节验收可以用一句话判断：能否从朴素版本推导到改进版本，并能说清新增字段、新增状态或新增实验到底保护了什么。如果只能给最终代码，却解释不了这条推导链，说明机制还没有真正学会。

\topic{工程边界：超时}

先看一个具体问题：一次 I/O 请求慢到什么程度算失败。在读取输入文件、写 shuffle block、提交 checkpoint中，这个问题会沿着输入、执行、同步、I/O 或提交边界继续传播，不会停留在一个局部函数里。

朴素写法通常是：无限等待完成。它吸引人的地方是短、直观、容易在小样本上通过测试；危险也在这里，因为小样本会替它隐藏资源、调度、数据分布和失败路径。

放大以后会出现的信号是：上游资源被长期占用，调用链预算失控。这个信号未必是崩溃，也可能是吞吐平台期、p99 抖动、结果偶尔变化、队列持续变长、重启后状态不一致，或者某个分片长期拖尾。

此时引入超时，不是为了增加术语，而是为了建立一个可检查模型：timeout 是资源预算的一部分，需要和重试、取消、幂等配合。模型必须同时解释正确性和成本。只解释为什么快，不解释为什么可信，不能进入贯穿项目；只解释为什么正确，却完全无视缓存、调度、I/O 或网络，也不符合第二册目标。

把模型写成不变量时，要问四个问题：哪些状态只能由一个拥有者修改，哪些结果允许重试，哪些数据可以批量搬运，哪些错误必须外显。围绕超时的代码审查，就从这四个问题开始，而不是从实现看起来是否高级开始。

观察方式应围绕假设设计：注入慢 I/O，观察超时后状态清理。实验要固定输入和环境，保留 reference，一次只改变一个变量。如果某项工具在当前手机或 Linux 环境没有权限，就写出预期命令、缺失事件和结论限制。

最后要写边界：底层操作可能超时后仍完成，不能假设自动撤销。边界不是附录里的保守声明，而是正文的一部分。读者应该知道什么时候使用它，什么时候退回简单方案，什么时候需要换一个尺度重新测。

本节验收可以用一句话判断：能否从朴素版本推导到改进版本，并能说清新增字段、新增状态或新增实验到底保护了什么。如果只能给最终代码，却解释不了这条推导链，说明机制还没有真正学会。

\topic{工程边界：可靠写入}

先看一个具体问题：write 返回成功是否等于 checkpoint 可恢复。在读取输入文件、写 shuffle block、提交 checkpoint中，这个问题会沿着输入、执行、同步、I/O 或提交边界继续传播，不会停留在一个局部函数里。

朴素写法通常是：写文件后立即认为提交。它吸引人的地方是短、直观、容易在小样本上通过测试；危险也在这里，因为小样本会替它隐藏资源、调度、数据分布和失败路径。

放大以后会出现的信号是：数据可能只在 page cache，目录项也可能未持久化。这个信号未必是崩溃，也可能是吞吐平台期、p99 抖动、结果偶尔变化、队列持续变长、重启后状态不一致，或者某个分片长期拖尾。

此时引入可靠写入，不是为了增加术语，而是为了建立一个可检查模型：可靠写入需要临时文件、fsync、rename 和必要的目录同步。模型必须同时解释正确性和成本。只解释为什么快，不解释为什么可信，不能进入贯穿项目；只解释为什么正确，却完全无视缓存、调度、I/O 或网络，也不符合第二册目标。

把模型写成不变量时，要问四个问题：哪些状态只能由一个拥有者修改，哪些结果允许重试，哪些数据可以批量搬运，哪些错误必须外显。围绕可靠写入的代码审查，就从这四个问题开始，而不是从实现看起来是否高级开始。

观察方式应围绕假设设计：用崩溃点矩阵分析提交前后。实验要固定输入和环境，保留 reference，一次只改变一个变量。如果某项工具在当前手机或 Linux 环境没有权限，就写出预期命令、缺失事件和结论限制。

最后要写边界：不同文件系统和设备语义有差异，教学项目要写边界。边界不是附录里的保守声明，而是正文的一部分。读者应该知道什么时候使用它，什么时候退回简单方案，什么时候需要换一个尺度重新测。

本节验收可以用一句话判断：能否从朴素版本推导到改进版本，并能说清新增字段、新增状态或新增实验到底保护了什么。如果只能给最终代码，却解释不了这条推导链，说明机制还没有真正学会。

\topic{工程边界：完成队列}

先看一个具体问题：多个 I/O 完成事件如何交给计算线程。在读取输入文件、写 shuffle block、提交 checkpoint中，这个问题会沿着输入、执行、同步、I/O 或提交边界继续传播，不会停留在一个局部函数里。

朴素写法通常是：完成回调直接做大量计算。它吸引人的地方是短、直观、容易在小样本上通过测试；危险也在这里，因为小样本会替它隐藏资源、调度、数据分布和失败路径。

放大以后会出现的信号是：回调线程被阻塞，完成队列堆积。这个信号未必是崩溃，也可能是吞吐平台期、p99 抖动、结果偶尔变化、队列持续变长、重启后状态不一致，或者某个分片长期拖尾。

此时引入完成队列，不是为了增加术语，而是为了建立一个可检查模型：完成队列把 I/O 事件转换为运行时任务。模型必须同时解释正确性和成本。只解释为什么快，不解释为什么可信，不能进入贯穿项目；只解释为什么正确，却完全无视缓存、调度、I/O 或网络，也不符合第二册目标。

把模型写成不变量时，要问四个问题：哪些状态只能由一个拥有者修改，哪些结果允许重试，哪些数据可以批量搬运，哪些错误必须外显。围绕完成队列的代码审查，就从这四个问题开始，而不是从实现看起来是否高级开始。

观察方式应围绕假设设计：记录 completion 到任务开始的延迟。实验要固定输入和环境，保留 reference，一次只改变一个变量。如果某项工具在当前手机或 Linux 环境没有权限，就写出预期命令、缺失事件和结论限制。

最后要写边界：回调中应保持短小，复杂工作交给任务池。边界不是附录里的保守声明，而是正文的一部分。读者应该知道什么时候使用它，什么时候退回简单方案，什么时候需要换一个尺度重新测。

本节验收可以用一句话判断：能否从朴素版本推导到改进版本，并能说清新增字段、新增状态或新增实验到底保护了什么。如果只能给最终代码，却解释不了这条推导链，说明机制还没有真正学会。

\topic{工程边界：错误分类}

先看一个具体问题：I/O 错误都重试是否正确。在读取输入文件、写 shuffle block、提交 checkpoint中，这个问题会沿着输入、执行、同步、I/O 或提交边界继续传播，不会停留在一个局部函数里。

朴素写法通常是：失败就立即重试。它吸引人的地方是短、直观、容易在小样本上通过测试；危险也在这里，因为小样本会替它隐藏资源、调度、数据分布和失败路径。

放大以后会出现的信号是：权限、空间不足和数据损坏重试无效，忙碌和临时错误可以退避。这个信号未必是崩溃，也可能是吞吐平台期、p99 抖动、结果偶尔变化、队列持续变长、重启后状态不一致，或者某个分片长期拖尾。

此时引入错误分类，不是为了增加术语，而是为了建立一个可检查模型：错误分类决定重试、降级、失败和告警。模型必须同时解释正确性和成本。只解释为什么快，不解释为什么可信，不能进入贯穿项目；只解释为什么正确，却完全无视缓存、调度、I/O 或网络，也不符合第二册目标。

把模型写成不变量时，要问四个问题：哪些状态只能由一个拥有者修改，哪些结果允许重试，哪些数据可以批量搬运，哪些错误必须外显。围绕错误分类的代码审查，就从这四个问题开始，而不是从实现看起来是否高级开始。

观察方式应围绕假设设计：构造 busy、not found、permission、corrupt 四类错误。实验要固定输入和环境，保留 reference，一次只改变一个变量。如果某项工具在当前手机或 Linux 环境没有权限，就写出预期命令、缺失事件和结论限制。

最后要写边界：重试必须有预算和抖动，否则会放大故障。边界不是附录里的保守声明，而是正文的一部分。读者应该知道什么时候使用它，什么时候退回简单方案，什么时候需要换一个尺度重新测。

本节验收可以用一句话判断：能否从朴素版本推导到改进版本，并能说清新增字段、新增状态或新增实验到底保护了什么。如果只能给最终代码，却解释不了这条推导链，说明机制还没有真正学会。

\topic{案例推导：先有问题，再有答案}

本章案例不直接给最终实现。每个案例都按同一条教学链路展开：先描述输入和失败，再写朴素版本为什么自然，接着指出第一处证据，最后才引入改进。这样读者能学会推导，而不是只记住答案。

\topic{案例：慢写出背压}

场景是：reduce 产生输出比磁盘写入快。先把它缩到可以手工检查的规模，再扩到能触发真实性能或恢复问题的规模。小规模负责证明语义，大规模负责暴露成本，两者缺一不可。
朴素版本是：无界缓存输出 block。它必须保留在仓库里，名字可以直接标成 baseline 或 naive。保留它不是为了上线，而是为了让反例可复现。
改进版本是：有界写队列加生产者等待。改进说明要写清楚它改变了哪条路径：少搬了哪些字节，少争了哪条 cache line，少等了哪个队列，减少了哪些重复请求，或者把哪个半提交状态变成了可恢复状态。
验收重点是：观察内存峰值和尾延迟。报告里至少有一组反例。没有反例的案例很容易变成宣传文字，读者也无法判断机制的边界。

\topic{案例：批大小扫描}

场景是：同样总字节按不同 batch 写出。先把它缩到可以手工检查的规模，再扩到能触发真实性能或恢复问题的规模。小规模负责证明语义，大规模负责暴露成本，两者缺一不可。
朴素版本是：每条记录单独写。它必须保留在仓库里，名字可以直接标成 baseline 或 naive。保留它不是为了上线，而是为了让反例可复现。
改进版本是：按 4KiB、64KiB、1MiB 等批大小测试。改进说明要写清楚它改变了哪条路径：少搬了哪些字节，少争了哪条 cache line，少等了哪个队列，减少了哪些重复请求，或者把哪个半提交状态变成了可恢复状态。
验收重点是：报告吞吐和 p99 延迟。报告里至少有一组反例。没有反例的案例很容易变成宣传文字，读者也无法判断机制的边界。

\topic{案例：checkpoint 崩溃点}

场景是：写 checkpoint 时在不同步骤 kill 进程。先把它缩到可以手工检查的规模，再扩到能触发真实性能或恢复问题的规模。小规模负责证明语义，大规模负责暴露成本，两者缺一不可。
朴素版本是：直接覆盖 checkpoint。它必须保留在仓库里，名字可以直接标成 baseline 或 naive。保留它不是为了上线，而是为了让反例可复现。
改进版本是：临时文件加 manifest 提交。改进说明要写清楚它改变了哪条路径：少搬了哪些字节，少争了哪条 cache line，少等了哪个队列，减少了哪些重复请求，或者把哪个半提交状态变成了可恢复状态。
验收重点是：恢复只接受完整且校验通过的版本。报告里至少有一组反例。没有反例的案例很容易变成宣传文字，读者也无法判断机制的边界。

\topic{实验矩阵：让结论可复现}

用慢写入模拟器控制写带宽，扫描不同队列容量和 batch 大小；checkpoint 在 open、write、fsync、rename、manifest 各点杀进程。实验要把正常输入、边界输入、压力输入和错误输入分开。正常输入证明功能还在，边界输入证明不变量，压力输入暴露成本，错误输入验证恢复。

围绕阻塞 I/O、异步完成、背压、批处理，不要把所有变量一次改完。比如改变布局时固定线程数，改变线程数时固定输入分布，改变批大小时固定写入设备状态，改变重试策略时固定故障注入脚本。变量克制是性能报告可信的前提。

\begin{lstlisting}[numbers=none]
experiment matrix:
  input: normal / boundary / pressure / faulty
  version: reference / naive / improved / counterexample
  metrics: correctness / throughput / latency / resource / recovery
  evidence: command / environment / raw sample / conclusion limit
\end{lstlisting}

报告中的数字要能追溯到命令。只贴一张结果表不够；还要写 CPU 或设备环境、编译选项、输入生成方式、是否预热、是否绑定线程、是否清理 page cache、是否启用故障注入。做不到的项目要写成未验证风险。

如果改进版本更慢，也不是失败。负结果常常比正结果更能训练判断力。它可能说明瓶颈不在这里，输入规模太小，机制成本超过收益，或者朴素方案在当前边界下已经足够。工程能力包括知道什么时候不改。

\topic{Linux 源码阅读：从用户态现象倒推对象}

本章 Linux 阅读入口保持窄范围：\filepath{fs/read_write.c}、\filepath{mm/filemap.c}、\filepath{fs/sync.c}、\filepath{io_uring/io_uring.c}。阅读目的不是把内核子系统背下来，而是把一个用户态现象映射到内核对象和状态变化上。

阅读笔记建议固定四列：用户态操作、内核对象、状态变化、可观察指标。用户态操作可以是等待、缺页、写文件、调度、计时或消息收发；内核对象可以是 task、VMA、page、inode、wait queue、bio、socket buffer 或计数器。

源码里的数据结构要比函数名更重要。链表、红黑树、位图、引用计数、状态枚举、等待队列、页表和缓存结构，决定了机制如何落地。读者不需要一次读完整个内核，但要能解释一个最小现象的因果链。

如果当前 Termux 或设备内核无法直接对应源码路径，也要写清版本边界。源码阅读提供机制参考，运行时证据仍来自当前机器上的实验。把二者混为一谈，会让报告看起来权威，实际不可复现。

\topic{排查演练：把事故走完一遍}

排查从具体差异开始：内存上涨时先看写队列水位和生产者等待；恢复错误时检查是否直接读了临时文件而没有 manifest 校验。这一步要慢，不要跳。越早跳到熟悉的工具，越容易把系统带到错误方向。

排查记录应保留时间线。第一行写事故输入和版本，第二行写 reference 摘要，第三行写朴素版本摘要，第四行写第一次分歧，后面才写指标和修复。时间线能防止复盘变成事后故事。

\begin{lstlisting}[numbers=none]
debug timeline:
  1. freeze input and version
  2. run reference and save digest
  3. run naive or optimized version
  4. compare first divergent boundary
  5. inspect metrics for that boundary
  6. patch invariant and rerun matrix
\end{lstlisting}

修复时只改一个边界。若同时改数据布局、线程数、队列容量、批大小和提交协议，即使结果变好也无法解释为什么。高质量教材的案例要能被读者复做，而不是只让读者相信作者经验。

\topic{项目验收边界}

贯穿项目不是把每章代码堆到一起，而是让每章新增一个可证明的能力。能力必须有 reference、朴素实现、改进实现、实验脚本、指标输出和失败注入中的至少几项。

\begin{lstlisting}[numbers=none]
project checkpoints:
  1. 给写出队列设置容量
  2. 记录 I/O 提交和完成事件
  3. 实现临时文件加 manifest
  4. 加入慢 I/O 和失败注入
  5. 写背压传播图
\end{lstlisting}

每个检查点都要进入仓库。只在正文里说“应该实现”不够；读者需要看到代码边界、测试边界和报告边界。若某个检查点受当前设备限制，提交降级说明，而不是把它从质量标准里删除。

项目报告最后写三段：本章新增能力解决了什么失败；新增能力引入了什么成本；哪些结论能迁移，哪些结论必须在新机器、新输入或新故障条件下重测。这个复盘会把整本书连成一条主线。

\topic{深度补充：完整走读、反例和实验手册}

现在把本章再走深一层。前面的正文已经给出事故入口：reduce 生成输出比磁盘写得快，CPU 看起来很空，内存却持续上涨，最终 checkpoint 还留下半写文件。。这一节不再添加新的名词，而是把一次真实排查拆成可重复的学习动作。读者要能拿着这一节，在自己的机器上复现一个小版本，再把结论迁移到贯穿项目。

第一步仍然是固定任务：读取输入文件、写 shuffle block、提交 checkpoint。固定任务不是为了限制想象，而是为了让所有机制共享同一组输入、同一条状态链、同一个 reference 和同一套指标。如果每讲一个概念就换一个例子，读者很难看到概念之间的因果关系；如果一直围绕同一条链路推进，读者能看到一个失败怎样把下一个机制推出来。

第二步是拆出最小可复现实验。最小实验不等于玩具实验，它必须保留会触发本章核心失败的形状：有足够大的输入，有明确的边界条件，有可比较的 reference，有一个朴素版本，有一个改进版本，还有一个能证明边界的反例。缺少任意一项，实验都会变成看起来能跑、实际无法解释的演示。

\begin{lstlisting}[numbers=none]
minimum reproducible study:
  1. freeze input and reference digest
  2. run baseline with stage metrics
  3. inject one pressure or failure condition
  4. change one mechanism only
  5. compare correctness, cost, and recovery
  6. write the counterexample before the conclusion
\end{lstlisting}

第三步是把结果写成证据链，而不是写成心得。证据链的形状很固定：现象是什么，怀疑哪条边界，为什么朴素版本会在这里失败，改进版本保护了哪个不变量，哪些指标支持这个判断，哪些情况没有覆盖。这比直接写“使用某技术后性能提升”更慢，但它能防止误导，也能让后续章节复用结论。

\topic{单点分析：阻塞 I/O}

围绕阻塞 I/O，不要先记定义，先问它解决哪一个具体卡点：read 阻塞时 worker 到底在等什么。在本章的读取输入文件、写 shuffle block、提交 checkpoint里，这个卡点通常不会独立出现，而是和输入规模、数据分布、执行顺序、资源上限或故障注入一起出现。

朴素版本是：在计算线程中直接读写文件。它在小输入下往往能通过，因为小输入没有把隐藏假设压穿。但是教材不能停在“它错了”这句话上，要说明它为什么一开始会被写出来：它减少了状态，靠调用栈维持顺序，靠当前机器的资源余量掩盖等待，靠人工观察代替指标。

一旦规模或失败形状变化，就会出现：线程被 I/O 等待占住，计算资源和等待资源混在一起。排查这类信号时，先不要改实现。先把它翻译成不变量问题：是不是同一个状态被多个拥有者写，是否某个输出位置没有唯一归属，是否某个对象生命周期跨过了异步边界，是否某个提交点缺少身份和校验。

机制模型是：阻塞 I/O 简单但会把等待带入执行线程。把模型落到代码时，至少要能指出一个字段、一个状态转换、一个测试和一个指标。字段让语义进入结构体，状态转换让路径可审查，测试让反例固定，指标让性能判断不靠猜。

实验观察是：记录 blocked time、CPU 利用率和队列长度。这里要避免一个常见错误：看到工具输出后直接写结论。正确做法是先写假设，再写为什么这个指标能支持或否定假设。若指标和假设没有关系，就不要把它放进报告。

边界是：对小工具阻塞 I/O 可能足够，不必过度异步。边界要写在正文里，因为工程设计通常不是“会不会”，而是“在什么条件下值得”。一个机制可能在大输入、高并发、失败恢复中必要，却在单线程小工具中完全多余。能说清这个取舍，才说明真正理解了机制。

\topic{单点分析：异步完成}

围绕异步完成，不要先记定义，先问它解决哪一个具体卡点：提交请求后，结果何时回来，谁拥有 buffer。在本章的读取输入文件、写 shuffle block、提交 checkpoint里，这个卡点通常不会独立出现，而是和输入规模、数据分布、执行顺序、资源上限或故障注入一起出现。

朴素版本是：把异步调用当成立刻完成。它在小输入下往往能通过，因为小输入没有把隐藏假设压穿。但是教材不能停在“它错了”这句话上，要说明它为什么一开始会被写出来：它减少了状态，靠调用栈维持顺序，靠当前机器的资源余量掩盖等待，靠人工观察代替指标。

一旦规模或失败形状变化，就会出现：buffer 生命周期、错误返回和取消都不清楚。排查这类信号时，先不要改实现。先把它翻译成不变量问题：是不是同一个状态被多个拥有者写，是否某个输出位置没有唯一归属，是否某个对象生命周期跨过了异步边界，是否某个提交点缺少身份和校验。

机制模型是：异步 I/O 把提交和完成拆开，需要 completion 事件和所有权规则。把模型落到代码时，至少要能指出一个字段、一个状态转换、一个测试和一个指标。字段让语义进入结构体，状态转换让路径可审查，测试让反例固定，指标让性能判断不靠猜。

实验观察是：测试完成前释放 buffer 和取消场景。这里要避免一个常见错误：看到工具输出后直接写结论。正确做法是先写假设，再写为什么这个指标能支持或否定假设。若指标和假设没有关系，就不要把它放进报告。

边界是：异步增加状态机复杂度，必须有清晰封装。边界要写在正文里，因为工程设计通常不是“会不会”，而是“在什么条件下值得”。一个机制可能在大输入、高并发、失败恢复中必要，却在单线程小工具中完全多余。能说清这个取舍，才说明真正理解了机制。

\topic{单点分析：背压}

围绕背压，不要先记定义，先问它解决哪一个具体卡点：写出速度慢于计算速度时，上游应该怎么办。在本章的读取输入文件、写 shuffle block、提交 checkpoint里，这个卡点通常不会独立出现，而是和输入规模、数据分布、执行顺序、资源上限或故障注入一起出现。

朴素版本是：继续产生输出并排队。它在小输入下往往能通过，因为小输入没有把隐藏假设压穿。但是教材不能停在“它错了”这句话上，要说明它为什么一开始会被写出来：它减少了状态，靠调用栈维持顺序，靠当前机器的资源余量掩盖等待，靠人工观察代替指标。

一旦规模或失败形状变化，就会出现：内存增长、延迟增长，最终 OOM 或超时。排查这类信号时，先不要改实现。先把它翻译成不变量问题：是不是同一个状态被多个拥有者写，是否某个输出位置没有唯一归属，是否某个对象生命周期跨过了异步边界，是否某个提交点缺少身份和校验。

机制模型是：背压把下游容量反馈到上游，让生产速度受限制。把模型落到代码时，至少要能指出一个字段、一个状态转换、一个测试和一个指标。字段让语义进入结构体，状态转换让路径可审查，测试让反例固定，指标让性能判断不靠猜。

实验观察是：记录队列水位、生产者等待和丢弃策略。这里要避免一个常见错误：看到工具输出后直接写结论。正确做法是先写假设，再写为什么这个指标能支持或否定假设。若指标和假设没有关系，就不要把它放进报告。

边界是：背压策略是业务语义，不能只有技术默认。边界要写在正文里，因为工程设计通常不是“会不会”，而是“在什么条件下值得”。一个机制可能在大输入、高并发、失败恢复中必要，却在单线程小工具中完全多余。能说清这个取舍，才说明真正理解了机制。

\topic{单点分析：批处理}

围绕批处理，不要先记定义，先问它解决哪一个具体卡点：为什么每条记录一次 write 会慢。在本章的读取输入文件、写 shuffle block、提交 checkpoint里，这个卡点通常不会独立出现，而是和输入规模、数据分布、执行顺序、资源上限或故障注入一起出现。

朴素版本是：记录一来就写出。它在小输入下往往能通过，因为小输入没有把隐藏假设压穿。但是教材不能停在“它错了”这句话上，要说明它为什么一开始会被写出来：它减少了状态，靠调用栈维持顺序，靠当前机器的资源余量掩盖等待，靠人工观察代替指标。

一旦规模或失败形状变化，就会出现：系统调用、锁和设备提交成本被放大。排查这类信号时，先不要改实现。先把它翻译成不变量问题：是不是同一个状态被多个拥有者写，是否某个输出位置没有唯一归属，是否某个对象生命周期跨过了异步边界，是否某个提交点缺少身份和校验。

机制模型是：批处理让一次昂贵操作服务多条记录。把模型落到代码时，至少要能指出一个字段、一个状态转换、一个测试和一个指标。字段让语义进入结构体，状态转换让路径可审查，测试让反例固定，指标让性能判断不靠猜。

实验观察是：比较批大小、延迟分位数和吞吐。这里要避免一个常见错误：看到工具输出后直接写结论。正确做法是先写假设，再写为什么这个指标能支持或否定假设。若指标和假设没有关系，就不要把它放进报告。

边界是：批太大增加尾延迟和崩溃后重做成本。边界要写在正文里，因为工程设计通常不是“会不会”，而是“在什么条件下值得”。一个机制可能在大输入、高并发、失败恢复中必要，却在单线程小工具中完全多余。能说清这个取舍，才说明真正理解了机制。

\topic{单点分析：超时}

围绕超时，不要先记定义，先问它解决哪一个具体卡点：一次 I/O 请求慢到什么程度算失败。在本章的读取输入文件、写 shuffle block、提交 checkpoint里，这个卡点通常不会独立出现，而是和输入规模、数据分布、执行顺序、资源上限或故障注入一起出现。

朴素版本是：无限等待完成。它在小输入下往往能通过，因为小输入没有把隐藏假设压穿。但是教材不能停在“它错了”这句话上，要说明它为什么一开始会被写出来：它减少了状态，靠调用栈维持顺序，靠当前机器的资源余量掩盖等待，靠人工观察代替指标。

一旦规模或失败形状变化，就会出现：上游资源被长期占用，调用链预算失控。排查这类信号时，先不要改实现。先把它翻译成不变量问题：是不是同一个状态被多个拥有者写，是否某个输出位置没有唯一归属，是否某个对象生命周期跨过了异步边界，是否某个提交点缺少身份和校验。

机制模型是：timeout 是资源预算的一部分，需要和重试、取消、幂等配合。把模型落到代码时，至少要能指出一个字段、一个状态转换、一个测试和一个指标。字段让语义进入结构体，状态转换让路径可审查，测试让反例固定，指标让性能判断不靠猜。

实验观察是：注入慢 I/O，观察超时后状态清理。这里要避免一个常见错误：看到工具输出后直接写结论。正确做法是先写假设，再写为什么这个指标能支持或否定假设。若指标和假设没有关系，就不要把它放进报告。

边界是：底层操作可能超时后仍完成，不能假设自动撤销。边界要写在正文里，因为工程设计通常不是“会不会”，而是“在什么条件下值得”。一个机制可能在大输入、高并发、失败恢复中必要，却在单线程小工具中完全多余。能说清这个取舍，才说明真正理解了机制。

\topic{单点分析：可靠写入}

围绕可靠写入，不要先记定义，先问它解决哪一个具体卡点：write 返回成功是否等于 checkpoint 可恢复。在本章的读取输入文件、写 shuffle block、提交 checkpoint里，这个卡点通常不会独立出现，而是和输入规模、数据分布、执行顺序、资源上限或故障注入一起出现。

朴素版本是：写文件后立即认为提交。它在小输入下往往能通过，因为小输入没有把隐藏假设压穿。但是教材不能停在“它错了”这句话上，要说明它为什么一开始会被写出来：它减少了状态，靠调用栈维持顺序，靠当前机器的资源余量掩盖等待，靠人工观察代替指标。

一旦规模或失败形状变化，就会出现：数据可能只在 page cache，目录项也可能未持久化。排查这类信号时，先不要改实现。先把它翻译成不变量问题：是不是同一个状态被多个拥有者写，是否某个输出位置没有唯一归属，是否某个对象生命周期跨过了异步边界，是否某个提交点缺少身份和校验。

机制模型是：可靠写入需要临时文件、fsync、rename 和必要的目录同步。把模型落到代码时，至少要能指出一个字段、一个状态转换、一个测试和一个指标。字段让语义进入结构体，状态转换让路径可审查，测试让反例固定，指标让性能判断不靠猜。

实验观察是：用崩溃点矩阵分析提交前后。这里要避免一个常见错误：看到工具输出后直接写结论。正确做法是先写假设，再写为什么这个指标能支持或否定假设。若指标和假设没有关系，就不要把它放进报告。

边界是：不同文件系统和设备语义有差异，教学项目要写边界。边界要写在正文里，因为工程设计通常不是“会不会”，而是“在什么条件下值得”。一个机制可能在大输入、高并发、失败恢复中必要，却在单线程小工具中完全多余。能说清这个取舍，才说明真正理解了机制。

\topic{单点分析：完成队列}

围绕完成队列，不要先记定义，先问它解决哪一个具体卡点：多个 I/O 完成事件如何交给计算线程。在本章的读取输入文件、写 shuffle block、提交 checkpoint里，这个卡点通常不会独立出现，而是和输入规模、数据分布、执行顺序、资源上限或故障注入一起出现。

朴素版本是：完成回调直接做大量计算。它在小输入下往往能通过，因为小输入没有把隐藏假设压穿。但是教材不能停在“它错了”这句话上，要说明它为什么一开始会被写出来：它减少了状态，靠调用栈维持顺序，靠当前机器的资源余量掩盖等待，靠人工观察代替指标。

一旦规模或失败形状变化，就会出现：回调线程被阻塞，完成队列堆积。排查这类信号时，先不要改实现。先把它翻译成不变量问题：是不是同一个状态被多个拥有者写，是否某个输出位置没有唯一归属，是否某个对象生命周期跨过了异步边界，是否某个提交点缺少身份和校验。

机制模型是：完成队列把 I/O 事件转换为运行时任务。把模型落到代码时，至少要能指出一个字段、一个状态转换、一个测试和一个指标。字段让语义进入结构体，状态转换让路径可审查，测试让反例固定，指标让性能判断不靠猜。

实验观察是：记录 completion 到任务开始的延迟。这里要避免一个常见错误：看到工具输出后直接写结论。正确做法是先写假设，再写为什么这个指标能支持或否定假设。若指标和假设没有关系，就不要把它放进报告。

边界是：回调中应保持短小，复杂工作交给任务池。边界要写在正文里，因为工程设计通常不是“会不会”，而是“在什么条件下值得”。一个机制可能在大输入、高并发、失败恢复中必要，却在单线程小工具中完全多余。能说清这个取舍，才说明真正理解了机制。

\topic{单点分析：错误分类}

围绕错误分类，不要先记定义，先问它解决哪一个具体卡点：I/O 错误都重试是否正确。在本章的读取输入文件、写 shuffle block、提交 checkpoint里，这个卡点通常不会独立出现，而是和输入规模、数据分布、执行顺序、资源上限或故障注入一起出现。

朴素版本是：失败就立即重试。它在小输入下往往能通过，因为小输入没有把隐藏假设压穿。但是教材不能停在“它错了”这句话上，要说明它为什么一开始会被写出来：它减少了状态，靠调用栈维持顺序，靠当前机器的资源余量掩盖等待，靠人工观察代替指标。

一旦规模或失败形状变化，就会出现：权限、空间不足和数据损坏重试无效，忙碌和临时错误可以退避。排查这类信号时，先不要改实现。先把它翻译成不变量问题：是不是同一个状态被多个拥有者写，是否某个输出位置没有唯一归属，是否某个对象生命周期跨过了异步边界，是否某个提交点缺少身份和校验。

机制模型是：错误分类决定重试、降级、失败和告警。把模型落到代码时，至少要能指出一个字段、一个状态转换、一个测试和一个指标。字段让语义进入结构体，状态转换让路径可审查，测试让反例固定，指标让性能判断不靠猜。

实验观察是：构造 busy、not found、permission、corrupt 四类错误。这里要避免一个常见错误：看到工具输出后直接写结论。正确做法是先写假设，再写为什么这个指标能支持或否定假设。若指标和假设没有关系，就不要把它放进报告。

边界是：重试必须有预算和抖动，否则会放大故障。边界要写在正文里，因为工程设计通常不是“会不会”，而是“在什么条件下值得”。一个机制可能在大输入、高并发、失败恢复中必要，却在单线程小工具中完全多余。能说清这个取舍，才说明真正理解了机制。

\topic{案例三轮推导：慢写出背压}

第一轮只写 reference。输入场景是：reduce 产生输出比磁盘写入快。reference 的代码允许慢，甚至可以很笨，但语义必须清楚。它要说明哪些输入合法，哪些输入算错误，输出如何排序或比较，错误如何计数，重复运行是否得到同一摘要。

第二轮写朴素工程版本：无界缓存输出 block。这一版的价值是暴露直觉缺陷。写作时要诚实说明它为什么诱人：代码短，状态少，不需要额外结构，小样本容易通过。然后立刻给出反例，而不是直接跳到最终方案。

第三轮才写改进版本：有界写队列加生产者等待。改进不是装饰，而是针对第二轮反例的最小修复。如果修复引入了新的成本，也要同时写出来。例如局部合并会增加最终 merge，scan 会增加一轮 pass，manifest 会增加持久化开销，分片会增加元数据。

验收是：观察内存峰值和尾延迟。验收报告还要额外写一条“不适用条件”。例如输入太小时，复杂机制可能不值得；数据分布变化时，热点处理可能失效；设备不支持某些计数器时，性能结论只能保留为假设。

\topic{案例三轮推导：批大小扫描}

第一轮只写 reference。输入场景是：同样总字节按不同 batch 写出。reference 的代码允许慢，甚至可以很笨，但语义必须清楚。它要说明哪些输入合法，哪些输入算错误，输出如何排序或比较，错误如何计数，重复运行是否得到同一摘要。

第二轮写朴素工程版本：每条记录单独写。这一版的价值是暴露直觉缺陷。写作时要诚实说明它为什么诱人：代码短，状态少，不需要额外结构，小样本容易通过。然后立刻给出反例，而不是直接跳到最终方案。

第三轮才写改进版本：按 4KiB、64KiB、1MiB 等批大小测试。改进不是装饰，而是针对第二轮反例的最小修复。如果修复引入了新的成本，也要同时写出来。例如局部合并会增加最终 merge，scan 会增加一轮 pass，manifest 会增加持久化开销，分片会增加元数据。

验收是：报告吞吐和 p99 延迟。验收报告还要额外写一条“不适用条件”。例如输入太小时，复杂机制可能不值得；数据分布变化时，热点处理可能失效；设备不支持某些计数器时，性能结论只能保留为假设。

\topic{案例三轮推导：checkpoint 崩溃点}

第一轮只写 reference。输入场景是：写 checkpoint 时在不同步骤 kill 进程。reference 的代码允许慢，甚至可以很笨，但语义必须清楚。它要说明哪些输入合法，哪些输入算错误，输出如何排序或比较，错误如何计数，重复运行是否得到同一摘要。

第二轮写朴素工程版本：直接覆盖 checkpoint。这一版的价值是暴露直觉缺陷。写作时要诚实说明它为什么诱人：代码短，状态少，不需要额外结构，小样本容易通过。然后立刻给出反例，而不是直接跳到最终方案。

第三轮才写改进版本：临时文件加 manifest 提交。改进不是装饰，而是针对第二轮反例的最小修复。如果修复引入了新的成本，也要同时写出来。例如局部合并会增加最终 merge，scan 会增加一轮 pass，manifest 会增加持久化开销，分片会增加元数据。

验收是：恢复只接受完整且校验通过的版本。验收报告还要额外写一条“不适用条件”。例如输入太小时，复杂机制可能不值得；数据分布变化时，热点处理可能失效；设备不支持某些计数器时，性能结论只能保留为假设。

\topic{实验手册：从命令到结论}

本章实验覆盖的机制包括：阻塞 I/O、异步完成、背压、批处理、超时、可靠写入、完成队列、错误分类。实验不是把这些名字逐个跑一遍，而是从同一条主线中抽取变量。先固定数据，再固定环境，再选择一个变量。变量可以是线程数、批大小、布局、分区函数、队列容量、故障注入点或重试预算，但一次只动一个。

实验输入建议分四类。正常输入用于确认功能没有被改坏；边界输入用于打到 off-by-one、空输入、满队列、零命中、重复消息等路径；压力输入用于暴露缓存、TLB、锁、队列、I/O 或网络成本；错误输入用于检查恢复、取消、幂等和错误分类。

每次实验至少记录五类结果。第一是 correctness digest，证明结果可信；第二是 throughput 或 elapsed time，说明端到端效果；第三是资源指标，例如内存峰值、队列水位、文件数量或 in-flight 数；第四是等待或失败指标，例如锁等待、超时、重试、late reply；第五是环境信息，例如 CPU、内核、编译器、输入规模和运行命令。

\begin{lstlisting}[numbers=none]
report skeleton:
  problem: one concrete symptom
  hypothesis: one boundary or cost source
  baseline: reference and naive version
  change: one mechanism only
  data: raw samples and digest
  result: what changed and what did not
  limit: unverified environment or counterexample
\end{lstlisting}

结论必须包含“没有变化”的部分。如果吞吐提高但内存峰值上升，要写；如果平均延迟降低但 p99 变差，要写；如果正确性通过但恢复路径未验证，要写。高质量教材不是让所有结果看起来漂亮，而是让读者学会在证据边界内说话。

\topic{Linux 源码映射：对象、状态和等待}

本章推荐源码入口是 \filepath{fs/read_write.c}、\filepath{mm/filemap.c}、\filepath{fs/sync.c}、\filepath{io_uring/io_uring.c}。阅读时不要展开成百科式笔记，而要围绕一个最小用户态现象。比如线程为什么睡眠，缺页为什么发生，写文件为什么没有立刻持久化，计数器为什么需要权限，队列为什么会唤醒。

源码阅读的核心是对象映射。用户态的线程，在内核里要找到 task 和调度状态；用户态的内存访问，要找到 VMA、页表项和 page；用户态的文件写入，要找到 file、inode、page cache 和 writeback；用户态的等待，要找到 wait queue、futex 或完成事件。

读内核代码时要特别注意状态字段如何变化。一个对象从 runnable 到 sleeping，从 clean page 到 dirty page，从 unlocked 到 contended，从 pending I/O 到 completed I/O，这些状态变化比函数名更能解释现象。把状态变化和本章实验指标对应起来，源码阅读才不是资料堆砌。

当前设备和源码版本可能不完全一致。报告要写明内核版本、阅读的源码版本、哪些结论来自源码机制、哪些结论来自本机实验。这能避免把源码阅读当作不可质疑的证明，也能训练读者区分机制理解和运行时证据。

\topic{贯穿项目验收：把能力写进系统}

本章进入贯穿项目时，不能只留下文字。至少要有一个目录、一个输入生成脚本或固定样本、一个 reference、一个朴素版本、一个改进版本、一个实验报告。代码可以小，但边界必须完整。

\begin{lstlisting}[numbers=none]
acceptance checklist:
  1. 给写出队列设置容量
  2. 记录 I/O 提交和完成事件
  3. 实现临时文件加 manifest
  4. 加入慢 I/O 和失败注入
  5. 写背压传播图
\end{lstlisting}

每个检查点都要对应一个失败路径。比如某个检查点是队列容量，就要能触发队列满；某个检查点是 manifest，就要能触发半写文件；某个检查点是 route epoch，就要能触发旧请求。没有失败注入的能力很容易只是正常路径演示。

验收时还要写回退策略。若改进版本复杂度太高、收益不足或设备环境不支持，项目应该能退回朴素但正确的版本。能回退说明边界清楚，不能回退往往说明机制已经和业务逻辑缠在一起。

本章结束前，读者应写一页复盘：本章新增了什么能力，解决了什么事故，引入了什么成本，下一章会复用哪个边界。这页复盘能让整本书保持连续，而不是每章都从零开始。

\topic{本节收束：回到最初事故}

最后再回到开头的事故：reduce 生成输出比磁盘写得快，CPU 看起来很空，内存却持续上涨，最终 checkpoint 还留下半写文件。。如果现在重新排查，读者不应只说“可能是某机制的问题”，而应能写出一条可执行路线：固定输入，运行 reference，复现朴素失败，选择一个机制，改一个边界，采集指标，写出反例和结论限制。

这种路线就是本章真正要交付的能力。知识点会随着硬件、内核和框架变化而更新，但从事故出发、沿边界推导、用证据收束的习惯不会过时。

\topic{工程补充：实现走读与反例库}

把本章写成教材，最终要能落到一段能维护的工程实现。围绕读取输入文件、写 shuffle block、提交 checkpoint，实现顺序建议从慢写出背压开始，因为它最容易把本章主失败暴露出来。先做 reference，再做朴素版本，最后做改进版本。reference 不追求速度，但它定义输入合同和输出摘要；朴素版本保留错误直觉；改进版本只改一个边界。

第一天实现不要碰太多机制。只写固定输入、固定输出和最小指标。输入要能触发这个事故：reduce 生成输出比磁盘写得快，CPU 看起来很空，内存却持续上涨，最终 checkpoint 还留下半写文件。。输出摘要不要只用一个总数，要至少包含记录数、错误数、分片数、关键输出 checksum 和阶段状态。这样后面每次改布局、线程、队列或协议，都能回到同一个摘要比较。

第二天再引入批大小扫描。这一类案例通常负责暴露资源或调度问题。写实现时要特别注意观测点放在哪里：放在热循环里会改变成本，放在边界处又可能看不到细节。教学项目可以先用较粗粒度的阶段指标，等路径稳定后再增加更细的 trace。

第三天引入checkpoint 崩溃点。这一类案例通常负责恢复、尾部、长尾或提交边界。此时不要让改进版本直接覆盖朴素版本。两个版本应能在同一输入上并排运行，并输出同一套摘要。并排运行可以让读者看到差异来自机制，而不是来自输入变化或测量变化。

实现走读还要保持模块边界。输入生成、reference、朴素版本、改进版本、指标采集、报告输出最好分成独立函数或独立小文件。不要为了写得快，把实验、业务逻辑和日志打印都塞进一个函数。这样的代码短期能跑，长期无法解释，也无法被后续章节复用。

\topic{状态走读：哪些事实必须显式记录}

读取输入文件、写 shuffle block、提交 checkpoint里最容易出错的地方，是把状态留在脑子里或调用栈里。只要出现线程、异步、重试、I/O 或分布式边界，调用栈就不再足够。教材要逼读者把事实写成字段：当前处理到哪里，谁拥有数据，哪次尝试有效，哪个输出已提交，哪个错误可重试，哪个指标能解释等待。

围绕阻塞 I/O，要记录的是语义事实。它回答“结果为什么可信”。如果这个事实没有字段承载，测试只能在结果出错后才发现问题；如果字段存在，状态机就能在错误路径上提前拒绝非法转换。

围绕异步完成，要记录的是成本事实。它回答“为什么慢或为什么不扩展”。成本事实不一定进入业务结构，但必须进入实验报告。例如队列水位、批大小、任务等待、cache miss、重试次数、partition 大小，这些数字决定了优化方向。

围绕背压，要记录的是边界事实。它回答“什么时候可以进入下一阶段”。边界事实尤其适合写成枚举状态，而不是一组布尔值。多个布尔值会产生 impossible state，枚举状态能让代码审查直接检查允许的转移。

状态走读的输出可以是一张简单表：状态名、拥有者、进入条件、退出条件、持久化要求、观测指标。读者不需要一开始画出完美状态机，但必须知道哪些状态不能靠注释维持。凡是会跨线程、跨文件、跨消息或跨重启的状态，都应优先显式化。

\begin{lstlisting}[numbers=none]
state review table:
  state name
  owner
  allowed incoming events
  allowed outgoing events
  persisted or transient
  metric that proves progress
\end{lstlisting}

\topic{反例库：防止机制变成口号}

反例库的作用，是防止读者把本章机制学成一句正确但空泛的话。每个反例都要小到可以复现，大到能代表真实系统的一类失败。反例不一定导致崩溃，很多时候它只表现为吞吐不上升、输出顺序不稳定、恢复后重复、队列增长或长尾放大。

反例：如果只按阻塞 I/O的直觉写代码，而不检查“read 阻塞时 worker 到底在等什么”，就会落回朴素路径“在计算线程中直接读写文件”。触发条件可以设计成：扩大输入、改变分布、插入等待、打乱顺序、制造重复或提前关闭。预期失败信号是：线程被 I/O 等待占住，计算资源和等待资源混在一起。这个反例应成为测试或实验的一部分，而不是只写在正文里。

反例：如果只按异步完成的直觉写代码，而不检查“提交请求后，结果何时回来，谁拥有 buffer”，就会落回朴素路径“把异步调用当成立刻完成”。触发条件可以设计成：扩大输入、改变分布、插入等待、打乱顺序、制造重复或提前关闭。预期失败信号是：buffer 生命周期、错误返回和取消都不清楚。这个反例应成为测试或实验的一部分，而不是只写在正文里。

反例：如果只按背压的直觉写代码，而不检查“写出速度慢于计算速度时，上游应该怎么办”，就会落回朴素路径“继续产生输出并排队”。触发条件可以设计成：扩大输入、改变分布、插入等待、打乱顺序、制造重复或提前关闭。预期失败信号是：内存增长、延迟增长，最终 OOM 或超时。这个反例应成为测试或实验的一部分，而不是只写在正文里。

反例：如果只按批处理的直觉写代码，而不检查“为什么每条记录一次 write 会慢”，就会落回朴素路径“记录一来就写出”。触发条件可以设计成：扩大输入、改变分布、插入等待、打乱顺序、制造重复或提前关闭。预期失败信号是：系统调用、锁和设备提交成本被放大。这个反例应成为测试或实验的一部分，而不是只写在正文里。

反例：如果只按超时的直觉写代码，而不检查“一次 I/O 请求慢到什么程度算失败”，就会落回朴素路径“无限等待完成”。触发条件可以设计成：扩大输入、改变分布、插入等待、打乱顺序、制造重复或提前关闭。预期失败信号是：上游资源被长期占用，调用链预算失控。这个反例应成为测试或实验的一部分，而不是只写在正文里。

反例：如果只按可靠写入的直觉写代码，而不检查“write 返回成功是否等于 checkpoint 可恢复”，就会落回朴素路径“写文件后立即认为提交”。触发条件可以设计成：扩大输入、改变分布、插入等待、打乱顺序、制造重复或提前关闭。预期失败信号是：数据可能只在 page cache，目录项也可能未持久化。这个反例应成为测试或实验的一部分，而不是只写在正文里。

反例库还要写出修复后如何证明不再发生。证明不能只靠“这次没有复现”。更好的方式是把反例输入固定下来，把状态摘要固定下来，把失败路径上的关键指标固定下来。以后重构时，这些反例就是回归测试。

\topic{实验协议：同一脚本跑出四类结论}

本章实验脚本要围绕一句话展开：用慢写入模拟器控制写带宽，扫描不同队列容量和 batch 大小；checkpoint 在 open、write、fsync、rename、manifest 各点杀进程。。脚本至少要支持四种模式：correctness、pressure、fault、report。correctness 模式跑小输入和边界输入，pressure 模式跑大输入或高并发，fault 模式启用故障注入，report 模式把原始样本整理成表。

correctness 模式不应该被性能优化绕过。它要在每个版本后运行，覆盖空输入、单元素、非整齐长度、重复 key、非法记录、关闭边界和恢复边界。如果某章没有这些具体路径，也要选择与本章机制等价的边界输入。

pressure 模式要把瓶颈逼出来。输入太小，很多机制看起来都没有意义；输入过大，又可能被外部环境噪声掩盖。合适的做法是用几档规模扫描，找到性能曲线的拐点，再围绕拐点解释。只报告最大规模，往往看不到瓶颈迁移过程。

fault 模式要把错误变成输入。错误不应依赖人工按 Ctrl-C 或碰运气，而应在脚本里用固定种子、固定时刻、固定任务或固定文件触发。例如在写 manifest 前杀进程，在某个消息上丢响应，在某个 partition 上制造热点，在某个 worker 上注入慢 I/O。

report 模式要保留原始样本，而不是只保留平均值。报告可以生成 CSV、纯文本摘要或 LaTeX 表格，但至少要能重新计算平均、分位数、最大值和错误计数。如果读者只能看到最终一句结论，就无法判断结论是否可靠。

\begin{lstlisting}[numbers=none]
lab modes:
  correctness: small and boundary inputs
  pressure: scale, concurrency, and distribution sweep
  fault: deterministic failure injection
  report: raw samples, digest, and conclusion limits
\end{lstlisting}

\topic{源码审查：从实现反推本章原则}

源码审查时不要先看函数数量，也不要先看是否用了某个高级技巧。先从数据入口开始，一路跟到提交或输出。每经过一个边界，就问这个边界有没有字段、状态、测试和指标。这个方法比逐行读代码更有效，因为它直接对应系统失败路径。

以本章为例，最小走查路线是：内存上涨时先看写队列水位和生产者等待；恢复错误时检查是否直接读了临时文件而没有 manifest 校验。。这条路线里，每个动作都应该能在源码中找到对应位置。如果找不到，说明代码把语义藏在约定里；如果能找到但没有测试，说明边界没有被固定；如果有测试但没有指标，说明性能和恢复结论仍然靠猜。

审查还要看错误路径。很多项目的正常路径清晰，错误路径却散在多个 return、异常、回调和清理函数里。教材里的示例必须展示错误如何返回、资源如何释放、等待者如何唤醒、临时输出如何清理、重复请求如何去重。否则读者学到的只是 happy path。

最后看命名。状态字段、计数器、日志事件和文件名都应体现语义身份。一个叫 data、flag、tmp、done 的字段，在小程序里能理解，在系统里会变成歧义源。更好的名字应包含对象、阶段和语义，例如 \code{current_attempt}、\code{committed_checksum}、\code{producer_wait_ns}、\code{route_epoch}。

\topic{迁移说明：从本章到下一章}

本章内容不能学完就断掉。每个机制都要问它会在后续哪里复用。单核里的依赖链会迁移到 SIMD 和并行归约；cache line 所有权会迁移到锁、原子和分片计数；有界队列会迁移到 I/O 背压和分布式流控；manifest 会迁移到 shuffle 和 checkpoint。

对于读取输入文件、写 shuffle block、提交 checkpoint，可迁移的是语义边界和实验方法，不可直接迁移的是具体数字。吞吐、延迟、批大小、线程数、队列容量、partition 数量，都要在新输入和新机器上重新测。这也是为什么本书反复要求记录环境和原始样本。

迁移时还要防止过度抽象。看到两个章节都有“队列”，不代表它们可以共用同一个类；看到两个章节都有“状态机”，不代表可以写一个万能状态机框架。抽象只有在减少真实重复、保护共同不变量、并且不会隐藏关键成本时才值得引入。

读者可以在每章末尾写三句话。第一，本章最重要的不变量是什么。第二，本章最可能迁移到后续哪一章。第三，本章最不能照搬的性能数字是什么。这三句话能帮助整本书形成连续推理，而不是把章节当作孤立文章。

\topic{章末审查：验收题、反例输入和提交标准}

本章最后再做一次审查。审查不是重复正文，而是把读取输入文件、写 shuffle block、提交 checkpoint压成可检查的交付物。最小验收题应覆盖这些案例：慢写出背压、批大小扫描、checkpoint 崩溃点。如果读者只能描述概念，却不能为这些案例写出 reference、朴素版本、反例和改进版本，就还没有达到本章要求。

第一类验收题检查语义。给定一组固定输入，要求写出串行 reference 和结果摘要。摘要要能暴露内存上涨时先看写队列水位和生产者等待；恢复错误时检查是否直接读了临时文件而没有 manifest 校验。。答案必须说明哪些字段参与比较，哪些错误要计数，哪些状态可以忽略，哪些状态必须持久化或记录。

第二类验收题检查机制推导。要求从阻塞 I/O、异步完成、背压、批处理、超时中选择一个机制，先写朴素方案为什么自然，再写它在本章主线中如何失败，最后写一个只改变单一边界的改进。答案不能直接给最终方案。没有朴素版本，就没有推导；没有失败证据，就没有必要性；没有反例，就没有边界。

第三类验收题检查实验设计。要求给出正常输入、边界输入、压力输入和错误输入四组样本。每组样本都要说明目的：验证功能、打到边界、暴露成本、检查恢复。如果一个样本只能让程序跑完，却不能支持或否定某个假设，就不应进入实验矩阵。

\begin{lstlisting}[numbers=none]
chapter review checklist:
  reference digest is stable
  naive version is preserved
  counterexample input is committed
  improved version changes one boundary
  metrics explain the chosen hypothesis
  report states unverified limits
\end{lstlisting}

\topic{反例输入怎么设计}

反例输入不要追求稀奇，而要追求直击本章边界。边界输入通常比大输入更有价值，因为它能把 off-by-one、空队列、重复请求、旧 attempt、半写文件、错误路由、尾部元素、分片倾斜这些问题直接打出来。

围绕读取输入文件、写 shuffle block、提交 checkpoint，一个好的反例输入应满足三个条件。第一，人工可以解释预期结果；第二，朴素版本会稳定暴露问题；第三，改进版本通过后仍能说明机制边界。如果反例太随机，失败了也难解释；如果反例太简单，可能无法触发真实风险。

反例还要保留原始材料。输入文件、随机种子、故障注入配置、运行命令、环境摘要和输出 digest 都应进入仓库或报告。不要只在正文里描述“构造一个错误输入”。没有材料，下一次重构时无法回归。

\topic{项目提交前的自检}

提交本章项目代码前，先问六个问题。第一，reference 是否还在。第二，朴素版本是否还能跑。第三，改进版本是否只改变一个主要机制。第四，错误路径是否有测试。第五，指标是否能解释结论。第六，报告是否写明没有验证的硬件或系统边界。

如果这些问题有任何一个回答不清，就先不要把章节标为完成。系统书的质量来自可复查，而不是来自一次性写出大量文字。读者能复查，作者能复查，后续章节也能复用，才说明本章真正稳定。

最后回到本章事故：reduce 生成输出比磁盘写得快，CPU 看起来很空，内存却持续上涨，最终 checkpoint 还留下半写文件。。章末自检的意义，就是让读者面对同类事故时能不慌乱。先固定事实，再找第一个分歧，再选机制，再写实验。这个顺序会在整本第二册反复出现。

\

t

o

p

i

c

{

学

习

复

盘

：

把

本

章

能

力

带

到

下

一

章

}







离

开

本

章

前

，

读

者

应

能

把

读

取

输

入

文

件

、

写



s

h

u

f

f

l

e



b

l

o

c

k

、

提

交



c

h

e

c

k

p

o

i

n

t

讲

成

一

条

完

整

因

果

链

，

而

不

是

讲

成

一

组

概

念

。

起

点

是

事

故

：

r

e

d

u

c

e



生

成

输

出

比

磁

盘

写

得

快

，

C

P

U



看

起

来

很

空

，

内

存

却

持

续

上

涨

，

最

终



c

h

e

c

k

p

o

i

n

t



还

留

下

半

写

文

件

。

；

中

间

是

朴

素

方

案

、

失

败

路

径

、

状

态

字

段

和

实

验

矩

阵

；

终

点

是

能

进

入

贯

穿

项

目

的

代

码

边

界

。

如

果

复

述

时

只

能

说

出

术

语

，

却

说

不

清

第

一

个

分

歧

出

现

在

输

入

、

执

行

、

同

步

、

I

/

O

、

提

交

还

是

恢

复

，

就

应

该

回

到

本

章

案

例

重

新

走

一

遍

。







本

章

最

少

要

带

走

两

个

能

力

。

第

一

个

是

围

绕

阻

塞



I

/

O

建

立

语

义

边

界

：

哪

些

输

入

合

法

，

哪

些

状

态

可

见

，

哪

些

结

果

可

以

比

较

，

哪

些

错

误

必

须

外

显

。

第

二

个

是

围

绕

异

步

完

成

建

立

成

本

边

界

：

哪

些

字

节

被

搬

动

，

哪

些

线

程

在

等

待

，

哪

些

队

列

在

积

压

，

哪

些

重

试

或

失

败

放

大

了

工

作

量

。

语

义

边

界

让

结

果

可

信

，

成

本

边

界

让

优

化

有

方

向

，

两

者

缺

一

不

可

。







本

章

案

例

包

括

慢

写

出

背

压

、

批

大

小

扫

描

、

c

h

e

c

k

p

o

i

n

t



崩

溃

点

。

复

盘

时

不

要

只

看

最

终

改

进

版

本

，

还

要

回

看

朴

素

版

本

为

什

么

会

被

写

出

来

。

很

多

工

程

事

故

的

根

源

不

是

作

者

不

知

道

高

级

机

制

，

而

是

小

规

模

条

件

成

立

后

，

没

有

在

规

模

、

并

发

、

故

障

和

恢

复

条

件

变

化

时

重

新

审

查

边

界

。

教

材

反

复

保

留

朴

素

版

本

，

就

是

为

了

训

练

这

种

迁

移

审

查

能

力

。







把

本

章

带

到

下

一

章

时

，

只

迁

移

方

法

，

不

照

搬

数

字

。

r

e

f

e

r

e

n

c

e

、

反

例

、

单

变

量

实

验

、

状

态

表

、

指

标

和

结

论

边

界

可

以

迁

移

；

具

体

吞

吐

、

延

迟

、

线

程

数

、

批

大

小

、

队

列

容

量

、

分

片

数

量

都

必

须

重

新

测

。

这

条

原

则

能

避

免

读

者

把

某

台

机

器

上

的

经

验

写

成

普

遍

规

律

。







最

后

给

自

己

留

一

个

三

句

话

笔

记

：

本

章

保

护

的

核

心

不

变

量

是

什

么

；

本

章

新

增

的

成

本

或

复

杂

度

是

什

么

；

本

章

哪

个

实

验

最

值

得

在

后

续

章

节

复

用

。

这

三

句

话

比

背

诵

术

语

更

重

要

，

因

为

它

能

把

章

节

串

成

一

套

工

程

判

断

。

\topic{本章质量标准}

读完本章，读者应该能围绕读取输入文件、写 shuffle block、提交 checkpoint讲清完整推导：朴素方案为什么成立，哪条失败路径先出现，引入的机制保护了哪个不变量，实验如何证明或否定假设。

如果正文只列概念，没有具体输入、状态、时间线和反例，就不合格；如果只给最终代码，没有朴素版本和失败证据，也不合格；如果只给 benchmark 数字，没有语义 reference 和结论边界，同样不合格。

本章最终留下的是判断力：面对一个慢、错、卡死、重复提交或无法恢复的系统，先固定问题，再画边界，再写 reference，再做对照，再改机制。这个顺序比任何单个技巧更重要。

\section{本章落到贯穿项目}

Compute Systems Engine 在本章应实现一个本机异步模拟，不必立刻依赖真实网络。可以用一个生产者按固定速率生成批次，一个解析阶段消耗 CPU，一个写出阶段按可配置延迟模拟 I/O。阶段之间使用 bounded queue，记录队列长度、阻塞次数、丢弃次数、批大小、端到端延迟和每阶段吞吐。

然后加入三种背压策略：阻塞生产者、try submit 失败、批量合并。再加入超时和取消：请求超过延迟预算后标记过期，但底层完成事件仍要被消费。这个实验会直接服务分布式 shuffle：网络发送、reduce 端接收、重试和检查点都可以从这里演化出来。

\section{本章习题}

\begin{exercisebox}
习题一：设计一个三阶段流水线：读取、计算、写出。为每个阶段写出资源类型、队列容量、背压策略、指标和失败语义。

习题二：比较阻塞 submit、try submit 和 timed submit。分别说明调用者能感知哪些状态，适合哪些业务，不适合哪些业务。

习题三：设计批处理策略。给定数量阈值 N 和时间阈值 T，说明高流量和低流量下的行为，解释吞吐和尾延迟如何权衡。

习题四：写一份 \code{io_uring} 学习报告提纲。只要求解释提交队列、完成队列、in-flight 请求和取消语义，不要求完整实现。

习题五：构造一个重试风暴场景。说明没有退避和预算时系统怎样恶化，再设计带预算、退避和幂等 ID 的重试策略。
\end{exercisebox}

\topic{完成队列和回调线程}

异步系统通常有完成队列。I/O 提交后，完成事件进入 completion queue，某个线程负责取出事件并执行后续处理。这个线程不应该做重 CPU 工作，否则完成队列会积压，导致已经完成的 I/O 迟迟不被用户态处理。常见做法是完成线程只解析结果、更新状态，然后把 CPU 工作提交到计算池。

完成事件必须携带请求上下文 ID，而不是直接裸指针访问业务对象。若请求已经超时，上下文可能处于 canceled 状态；完成线程应安全地把结果丢弃或记录，而不是继续执行已过期业务逻辑。这里的状态机通常包括 submitted、completed、timed out、canceled、reaped。reaped 表示底层完成事件已经被消费，资源可以释放。

\topic{短读、短写和 partial completion}

I/O API 经常允许短读和短写。读取文件或 socket 时，返回的字节数可能小于请求字节数；写入也是如此。正确代码必须循环直到满足条件、遇到 EOF、出错或超时。若把一次 read/write 返回当成完整消息，就会在压力和边界情况下出错。

协议通常需要 framing：先读固定长度 header，再根据长度读 body；或者使用分隔符；或者使用固定大小消息。异步实现中，一个逻辑消息可能对应多次底层完成事件。请求状态机要记录已经读写多少字节、剩余多少字节、下一步做什么。I/O 编程的难点经常不在系统调用，而在状态机完整性。

\topic{限流器和令牌桶}

背压可以由队列容量触发，也可以由限流器主动控制。令牌桶是一种常见模型：系统按固定速率产生 token，请求需要消耗 token 才能发送；桶容量允许短时间突发；长期速率受 token 生成速度限制。它适合控制出站 RPC、日志写入、后台压缩和重试。

令牌桶的参数要有业务含义。生成速率对应可持续吞吐，桶容量对应允许突发，等待策略决定超出速率时阻塞还是失败。若重试也共用同一个令牌桶，就能防止重试风暴；若重试绕过限流，就会在故障时放大流量。

\topic{队列容量的选择}

队列容量不是越大越好。容量太小，短暂抖动就会阻塞上游，吞吐可能下降；容量太大，过载时延迟积压很深，失败反馈太晚，还可能耗尽内存。容量应根据服务时间、突发大小、延迟预算和内存预算选择。

一个实用方法是先估算：如果消费者每秒处理 10000 个任务，允许吸收 100 ms 突发，则队列容量至少约 1000；如果端到端延迟预算只有 50 ms，队列排队就不能无限增长。测量时要记录队列等待时间，而不只是队列长度。长度相同，处理速率不同，等待时间也不同。

\topic{优雅关闭和 drain}

I/O 系统关闭时要定义 drain 语义。立即关闭会丢弃队列中尚未处理的请求；drain 关闭会拒绝新请求，但继续处理已接受请求；带超时 drain 会在预算内尽力完成，超时后取消或丢弃。不同业务选择不同语义。日志系统可能允许丢低优先级日志，数据库提交不应随便丢。

线程池、bounded queue 和异步 I/O 都需要关闭协议。关闭不是析构时把标志设为 true 就结束，还要唤醒等待者、停止接受新请求、等待 running 请求、消费完成队列、释放资源并报告未完成请求。若完成队列没有 drain，底层稍后完成可能访问已经销毁的上下文。

\topic{本章项目验收}

Compute Systems Engine 的 I/O 模拟模块应通过这些验收：生产速度大于消费速度时队列不会无限增长；阻塞策略能降低生产速度；try 策略能报告失败；批处理策略能降低固定成本但不超过延迟预算；超时请求的完成事件仍被安全回收；关闭时不丢已接受且必须完成的任务；指标能显示队列长度、等待时间、批大小、超时数和丢弃数。

这些验收看起来像工程细节，但它们决定系统能否在过载和退出时保持可控。很多线上事故不是算法算错，而是队列无界、关闭不完整、重试没有预算、完成事件访问已释放对象。

\topic{文件 I/O 的提交语义}

文件写入有多层语义。应用把数据复制到用户态缓冲，只说明用户态内存有数据；调用 \code{write} 成功，通常说明数据进入内核 page cache 或设备队列，不一定持久化；调用 \code{fsync} 成功，才更接近“崩溃后可恢复”，但还要考虑文件系统、目录项和设备缓存。若用临时文件加 rename 提交，还要考虑 rename 后 fsync 目录，确保目录项持久化。

这对 checkpoint 和 manifest 非常关键。一个作业写完 manifest 后就宣布成功，但 manifest 没有 fsync，机器掉电后可能丢失。或者输出文件已经 rename，但目录项没持久化，恢复时看不到文件。教学项目可以在手机上简化持久化实验，但正文必须说明真实可靠提交需要明确持久化边界。

\begin{lstlisting}[numbers=none]
durable output shape:
  write temporary output
  fsync temporary output
  rename temporary output to final name
  fsync parent directory
  record manifest commit
\end{lstlisting}

不同文件系统和存储服务语义不同。对象存储可能没有 POSIX rename，NFS 语义也可能不同。本书不展开存储系统细节，但要求读者不要把“写了文件”误认为“已经安全提交”。

\topic{Page cache 对 benchmark 的影响}

文件读取 benchmark 很容易测到 page cache。第一次读大文件可能来自磁盘，第二次读同一文件可能来自内存。若报告没有说明冷 cache 还是热 cache，结果无法解释。清 page cache 需要权限，也会影响系统其他进程；普通用户可以通过读取不同文件、增大数据集或明确报告热缓存来处理。

写入 benchmark 也类似。写入 page cache 很快，后台写回稍后发生。若只测 write 返回时间，可能低估真实持久化成本；若测 fsync，成本会集中暴露。日志系统通常用批量 fsync 或 group commit，在吞吐和延迟之间取舍。第 17 和 18 章的 checkpoint 提交，需要继承这里的语义。

\topic{epoll 的 ready 语义}

\code{epoll} 通知的是 fd ready，而不是业务消息完成。Socket 可读表示读一次不会阻塞，但不保证完整应用消息已经到达；可写表示发送缓冲有空间，但不保证整个响应能一次写完。应用仍然要处理短读短写和消息 framing。事件循环只是等待机制，不替你实现协议状态机。

边沿触发和水平触发也有差异。水平触发下，只要状态仍 ready，事件会继续出现；边沿触发下，状态从 not ready 变 ready 时通知一次，应用必须读写到 EAGAIN，否则可能遗漏后续机会。边沿触发可以减少重复通知，但更容易写错。教学阶段建议先理解水平触发，再研究边沿触发。

\topic{网络缓冲和 backpressure}

网络发送通常先进入用户态缓冲、内核 socket buffer、网卡队列，再到网络。\code{send} 成功不表示对端应用已经处理，只表示本端接受了数据。若对端处理慢，接收窗口缩小，本端发送缓冲会填满，最终 \code{send} 阻塞或返回 EAGAIN。这是 TCP 层面的背压。

应用层还需要自己的背压。TCP 只能说明字节流压力，不知道业务优先级、请求幂等、队列容量和服务端 CPU。应用可以根据 reduce partition 队列长度返回 busy，可以根据请求 deadline 拒绝，可以按优先级丢弃低价值消息。网络背压和应用背压要配合，而不是互相替代。

\topic{零拷贝的生命周期问题}

零拷贝听起来总是好事，但它把缓冲区生命周期变长。普通发送可以把数据复制到内核缓冲后复用用户内存；零拷贝发送可能要求用户缓冲在发送完成前保持不变。若应用过早复用或释放缓冲，就会发送错误数据或崩溃。完成通知必须告诉应用何时可以回收缓冲。

这和异步 I/O 的请求上下文相同：提交出去的东西，在完成前不能随便销毁。对象池、引用计数、completion queue 和取消状态机都要配套。没有生命周期协议的零拷贝优化，会把一个性能问题变成正确性问题。

\topic{Group commit}

Group commit 把多个提交合并到一次持久化操作中。数据库日志、消息队列和检查点系统都常用这个思想。单个请求每次 fsync，延迟高且吞吐低；多个请求共享一次 fsync，可以大幅提高吞吐，但每个请求可能等待批次凑齐。数量阈值和时间阈值决定延迟和吞吐。

Group commit 还要处理失败。如果批次 fsync 失败，批内所有请求都失败还是部分失败？如果 fsync 成功但响应发送失败，客户端重试时如何返回旧结果？如果进程在 fsync 后、响应前崩溃，恢复时是否能根据日志判断请求已提交？这些问题直接连接分布式幂等表和提交点。

\topic{异步错误传播}

同步函数可以通过返回值或异常报告错误，异步系统的错误可能在提交之后很久才出现。提交成功只表示请求被接受，不表示操作成功。调用者需要一个 completion、future、回调或事件来接收最终结果。若接口只返回 accepted，却没有完成路径，错误会丢失。

异步错误还可能发生在取消之后。上层已经超时，底层返回失败或成功，完成路径仍要消费事件并更新状态。状态机应区分 user-visible result 和 internal cleanup。对用户来说请求已超时；对系统来说底层资源还没 reaped。混淆这两者会导致内存泄漏或 use-after-free。

\topic{队列的水位线}

除了满和空，队列还可以设置高水位线和低水位线。超过高水位线时上游减速或停止；降到低水位线时恢复。水位线可以避免在满和不满之间频繁抖动。网络协议、日志系统和消息队列都常用类似机制。

水位线选择要结合处理速率和突发。高水位线太接近容量，反馈太晚；太低，吞吐可能受限。低水位线太接近高水位线，系统频繁开关；太低，恢复太慢。实验可以画队列长度随时间变化，观察是否出现锯齿、振荡或持续增长。

\topic{优先级丢弃}

有些数据可以丢，有些不能丢。监控采样、调试日志、推荐候选可以在过载时丢弃或降级；支付请求、提交协议和检查点不能随意丢。队列满时，如果只用 FIFO，低价值任务可能占住空间，高价值任务被拒绝。优先级丢弃或多队列隔离可以改善这种情况。

但丢弃必须可观测。系统要记录丢弃数量、原因、优先级和影响范围。静默丢弃会让上层误以为系统正常。对于可丢数据，也要定义最大丢弃比例或降级策略，避免过载时完全失明。

\topic{背压传播链}

背压必须沿数据流传播到真正的源头。假设 write queue 满了，compute worker 被阻塞；compute worker 阻塞后 parse queue 满；parse queue 满后 read stage 停止读取；read stage 停止读取后 socket receive buffer 满；最终 TCP 窗口缩小，上游发送变慢。这是一条健康的传播链。若中间某处使用无限队列，背压就断了，内存会积压在那里。

分布式 shuffle 也一样。Reduce 队列满，应反馈给 map 端；map 端降低发送或合并更多；如果 map 端仍无限读取输入并缓存 shuffle，内存会爆。背压链条要在设计图里画出来，不能只在某个队列上写容量。

\topic{超时和 Deadline 的区别}

Timeout 是从当前操作开始计算的一段时间，deadline 是整个请求的截止时间。异步流水线更适合传 deadline，因为请求经过多个阶段，每个阶段都应知道剩余预算。若每个阶段都设置固定 timeout，端到端延迟可能超过用户预算。Deadline 还能帮助队列做决策：已经没有足够时间完成的任务，可以提前拒绝或降级。

Deadline 也要进入重试策略。一个请求剩余 5 ms，却启动一个预期 50 ms 的重试，通常没有意义。重试库应检查剩余预算、幂等性和重试次数。分布式系统中，deadline 还要考虑时钟差异；本机流水线中，单调时钟即可。

\topic{I/O 模拟器的参数}

本章项目的 I/O 模拟器应有可调参数：生产速率、消费速率、队列容量、批大小、批时间阈值、错误率、短写概率、超时预算、重试次数和是否允许丢弃。通过改变这些参数，读者可以观察队列从稳定到过载的过程。模拟器不是为了假装真实设备，而是为了训练状态机和背压思维。

每次实验要记录参数。否则“队列爆了”没有意义。记录后可以复现：在生产速率两倍于消费速率、队列容量 1024、批大小 64、错误率 1\% 时，阻塞策略如何表现，try 策略丢多少，批处理如何影响延迟。

\topic{I/O 章节和分布式章节的连接}

本章所有概念都会在分布式章节放大。短写对应消息分片，completion queue 对应 RPC 响应，deadline 对应跨服务预算，令牌桶对应重试限流，group commit 对应提交日志，manifest fsync 对应 checkpoint，背压链对应 shuffle 流控。学 I/O 不是为了记住某个系统调用，而是为了理解等待、完成、提交和过载这些通用结构。

\topic{端到端流水线案例}

以日志统计为例，端到端流水线可以这样设计。Read stage 从文件读取 batch，Parse stage 把文本转成列式热字段，Compute stage 做本地聚合，Shuffle stage 按 partition 发送 partial，Write stage 写 manifest 或最终输出。每两个 stage 之间都有有界队列。每个 batch 都有 batch id、输入 offset、记录数、deadline 和错误计数。

这个流水线里，任何阶段慢都会反映到队列。Parse 慢，read-to-parse 队列增长；Compute 慢，parse-to-compute 队列增长；Shuffle 慢，compute 输出缓冲增长；Write 慢，提交延迟上升。通过这些队列，系统可以定位瓶颈，而不是只看到总时间变长。队列也是背压传播路径，防止 read stage 无限读取文件占用内存。

\begin{lstlisting}[language=C++,caption={流水线 batch 元数据}]
struct PipelineBatch {
    std::uint64_t batch_id = 0;
    std::uint64_t begin_offset = 0;
    std::uint32_t record_count = 0;
    std::uint32_t error_count = 0;
};
\end{lstlisting}

Batch 元数据应和数据分开。热字段列可能很大，元数据用于调度、指标和恢复。后续分布式版本中，batch id 可以映射到 task id 或 shuffle block id。

\topic{可靠写入的状态机}

写输出可以分成 preparing、writing、flushing、committing、committed、failed。Preparing 创建临时路径和缓冲；writing 写数据；flushing 执行必要持久化；committing 发布 manifest 或 rename；committed 对外可见；failed 清理临时状态或保留诊断。把这些状态写清，恢复时才知道半成品如何处理。

如果进程在 writing 中崩溃，临时文件可能不完整，应丢弃或重写；在 flushing 后、committing 前崩溃，数据可能持久但未对外宣布，恢复时应根据 manifest 决定是否采用；在 committing 后响应丢失，重试应返回已提交。这个状态机和分布式 task attempt 提交完全一致，只是资源从网络输出变成文件。

\topic{读放大和写放大}

I/O 优化要关注读放大和写放大。读取一条记录可能需要读入整页、解压整个块或读取索引；写一条记录可能导致写一个 page、一个日志段或一个压缩块。若每条小记录单独写文件，元数据和 fsync 成本会非常高；批量写能降低写放大。若每次查询都读完整原始日志，只为了统计一个字段，读放大很高；列式存储和索引可以降低读放大。

读写放大和 cache line 放大是同一思想在存储层的版本。硬件按 cache line 搬，文件系统按页和块搬，网络按包和缓冲搬。系统设计要让每次搬动包含足够有用信息。

\topic{异步管线的内存预算}

异步流水线中，每个队列容量乘以 batch 大小，就是内存预算的一部分。若 read queue 容量 128，每个 batch 4 MiB，只这一段就可能占 512 MiB。多个阶段叠加后，内存峰值很快上升。队列容量不能只按任务数量设置，还要按字节设置。

因此实际系统常同时限制 item count 和 byte count。一个 batch 很大时，即使 item count 不高，也可能触发背压。Shuffle 阶段尤其需要按字节限流，因为不同 partition 的 block 大小差别很大。第 18 章的 map in-flight shuffle bytes 就来自这个原则。

\topic{完成回调的幂等}

异步完成回调可能被重复触发吗？理想 API 不应重复，但上层重试、取消和状态恢复可能让同一逻辑请求出现多个完成路径。回调内部应先检查请求状态，只有从 submitted 转到 completed 的那一次执行用户可见副作用。后到的完成、取消或超时只做清理或记录重复。

这可以用 CAS 状态转换表达。状态从 Submitted 到 Completed 成功的线程负责发布结果；若状态已经 TimedOut，完成线程只回收底层资源；若状态已经 Canceled，同样不执行业务回调。这个小状态机和分布式 exactly-once 提交同构。

\topic{I/O 失败的可恢复性}

不是所有 I/O 失败都可恢复。临时 EAGAIN、超时、连接重置可以重试；磁盘满、权限错误、校验失败、格式错误可能需要停止或人工处理。系统应把错误分为 retryable、fatal、corrupt、backpressure。Corrupt 错误尤其重要：如果读取 block 校验失败，盲目重试同一损坏文件可能无用，需要重新生成或从副本恢复。

错误分类进入 metrics。Retryable 错误升高可能说明下游过载；fatal 错误说明配置或权限；corrupt 错误说明存储或写入协议；backpressure 说明容量保护生效。没有分类，错误数只是噪声。

\topic{本章附加实验设计}

附加实验一：模拟 writer 每批写入后是否 fsync。比较只 write、每批 fsync、每 N 批 group fsync 的吞吐和延迟，并说明持久化语义差异。附加实验二：模拟队列按 item count 限制和按 byte count 限制，构造少量大 batch 和大量小 batch，观察哪种策略能控制内存。附加实验三：模拟 completion 在 timeout 后到达，验证资源能被 reaped 且业务回调不会重复执行。

这些实验会让 I/O 章和分布式章真正连接起来。可靠计算系统的难点不是一次成功写入，而是在超时、失败、重复和关闭时仍保持状态机正确。

\topic{案例：限速读取}

当下游处理慢时，读取阶段不应继续把文件读入内存。限速读取可以根据 parse queue 或 compute queue 的水位线决定是否暂停。若队列超过高水位线，read stage 暂停或减小 batch；低于低水位线后恢复。这样背压从 compute 传回 read。限速读取比读完再排队更稳定。

这个案例在本机文件上也成立。即使文件读取很快，如果解析和聚合慢，无限制读取会占用内存。真实网络输入更明显：停止读取 socket 会让 TCP 接收窗口缩小，把背压传给发送端。系统应尽量让背压接近源头，而不是在中间堆积。

\topic{案例：写出线程和计算线程分离}

Reduce 完成后要写输出。若 worker 在线程池里直接执行写文件和 fsync，它会占住计算 worker，其他 reduce 任务无法运行。更好的设计是计算 worker 生成 OutputBlock，提交给 write stage；write stage 有自己的并发度和队列。写慢时，write queue 背压 reduce；reduce 再背压 shuffle；最终 map 降速。

分离后也要处理提交结果。Write stage 成功后通知 driver commit，失败后通知 driver 重试或失败作业。Compute worker 不能假设提交立即成功。这个结构把 I/O 等待从计算池中剥离，但增加了异步状态机。

\topic{案例：超时后的输出}

一个写出请求可能在上层超时后仍然成功落盘。如果上层已经重试并提交了另一个 attempt，迟到的成功不能覆盖最终输出。解决方法是写出路径只写 attempt 临时文件，最终 manifest 由 driver 决定。迟到 attempt 即使写成功，也不会进入 manifest；清理线程稍后删除它。I/O 超时和分布式 attempt 去重在这里完全重合。

这个案例说明为什么写出路径不能直接写最终文件。直接写最终文件没有给 driver 选择获胜 attempt 的机会，也无法安全处理迟到完成。临时输出和 manifest 提交是可靠计算的基本结构。

\topic{端到端延迟分解}

流水线报告应把端到端延迟分解为队列等待和服务时间。一个 batch 的延迟可以拆成 read service、parse queue wait、parse service、compute queue wait、compute service、write queue wait、write service、commit wait。若只看总延迟，无法知道瓶颈在哪。若队列等待占比高，说明下游容量不足或背压；若 service 时间高，说明该阶段本身慢。

Trace 可以记录每个 batch 的阶段时间。聚合后可以看中位数和尾部。很多系统平均吞吐正常，但少数 batch 因为写出 fsync 或热点 key 尾部很长。端到端分解能帮助定位尾延迟。

\topic{I/O 背压验收清单}

本章最终验收应包含：队列有 item 和 byte 两种容量；生产者在高水位线暂停；低水位线恢复；写出使用临时文件和提交点；completion 迟到不会重复业务回调；超时请求最终 reaped；错误分类进入指标；关闭支持 drain 和立即取消两种策略；trace 能显示每个 batch 的等待和服务时间。达到这些要求，I/O 模块才足以支撑第 18 章的分布式计算。

\topic{I/O 章节总结}

I/O 系统的核心是等待和完成。同步阻塞简单但占线程，异步减少线程等待但增加状态机；批处理提高吞吐但增加延迟；有界队列暴露过载；deadline 和取消控制生命周期；临时文件和 manifest 定义提交。I/O 代码的质量不在于用了哪个 API，而在于是否能在短读、短写、超时、取消、关闭和失败时保持状态正确。

\topic{I/O 决策表}

设计 I/O 路径时，先问并发等待数量是否大，是否需要异步；数据是否必须持久化，是否需要 fsync 和 rename；是否允许丢弃，队列满时阻塞还是失败；是否有 deadline，超时后底层完成如何 reaped；是否需要批处理，数量阈值和时间阈值是多少；是否按 item 和 byte 双重限流；错误是否分类；关闭是 drain 还是立即取消。把这些答案写清，I/O 状态机才不会在边界条件下崩溃。

\topic{案例：检查点写入调度}

Checkpoint 写入可能和正常计算争用磁盘。若每个 stage 结束都立即写大 checkpoint，作业延迟会出现尖峰；若后台慢慢写，又要处理写未完成时进程崩溃。可以把 checkpoint 分成元数据和大数据：关键 manifest 小而同步提交，大块数据在 stage 输出时已经写好并校验。这样 checkpoint 提交只发布指针和摘要，减少关键路径 I/O。

这个案例展示了控制面和数据面分离在 I/O 层面的意义。真正必须同步持久化的是“哪些数据有效”的小元数据，大数据可以提前流式写入临时位置。第 18 章 manifest 正是这种结构。

\topic{案例：日志丢弃策略}

系统过载时，调试日志可以丢，但错误报告和提交日志不能丢。可以为日志设置优先级队列：critical 日志阻塞或落盘，debug 日志在队列满时丢弃并增加 drop counter。这样过载时系统仍保留关键证据，不会被低价值日志拖垮。

丢弃策略必须透明。Drop counter 要进入指标，报告中要说明哪些日志可能丢。静默丢弃会让故障排查失去证据。背压策略本质上是业务价值排序。

\topic{案例：异步写入和 backpressure 闭环}

假设 reduce worker 生成输出 block，write stage 写磁盘。如果 write queue 超过高水位，reduce worker 停止接收新的 shuffle block；reduce 队列满后 map worker 降低发送；map 输出缓冲满后 input reader 暂停读取。这个闭环把磁盘压力传到输入源。若任何一段使用无限队列，闭环断开，内存会在那段堆积。

设计背压时，要画出这条闭环，并标出每段的容量和策略。哪个阶段阻塞，哪个阶段返回 busy，哪个阶段丢弃，哪个阶段降级，都要明确。背压不是单个队列属性，而是整条数据流的协议。

\topic{案例：超时预算贯穿流水线}

一个 batch 从读取到提交可能有总预算。Read stage 用掉一部分，parse queue 等待用掉一部分，compute 用掉一部分，write fsync 又用掉一部分。如果每段只看自己的 timeout，端到端可能超预算。更好的方式是 batch 携带 deadline，每个阶段根据剩余时间决定继续、降级还是失败。

Deadline 也帮助背压。队列里等待太久的 batch，即使后续执行也可能超过预算，可以提前丢弃或标记失败，释放资源给仍有机会成功的任务。这个策略不适合必须完成的批处理，但适合交互请求和有时效性的日志分析。时间语义是调度的一部分。

\topic{本章扩展习题}

\begin{exercisebox}
习题六：设计 completion queue 的请求状态机。状态至少包含 submitted、completed、timed out、canceled、reaped。说明每个状态允许哪些转移。

习题七：写一个短读短写处理流程。要求支持 header/body 消息，记录已处理字节数，并说明错误和 EOF 如何处理。

习题八：为重试请求设计共享令牌桶。说明正常请求和重试请求是否共用 token，为什么。

习题九：设计 shutdown drain 策略。比较立即关闭、drain 关闭、带超时 drain 三种语义，以及它们适合的业务。
\end{exercisebox}

\begin{labbox}
实验：用 bounded queue 模拟生产者和消费者。让生产速度高于消费速度，观察队列大小和阻塞行为。解释这为什么是背压，而不是简单的性能问题。
\end{labbox}
